\documentclass[11pt]{article}
\usepackage[margin=1in]{geometry}
\usepackage{amsmath,amssymb,amsthm,mathtools,microtype,booktabs}
\usepackage[hidelinks]{hyperref}
\hypersetup{
 pdftitle={Talagrand's two-cost discrepancy problem: a prefix entropy criterion},
 pdfauthor={Michael Goldman, Francesco Grotto, Dario Trevisan}}
\newenvironment{aiuse}
  {\renewcommand{\abstractname}{AI use}\begin{abstract}}
  {\end{abstract}}
\newtheorem{theorem}{Theorem}[section]
\newtheorem{lemma}[theorem]{Lemma}
\newtheorem{proposition}[theorem]{Proposition}
\newtheorem{corollary}[theorem]{Corollary}
\newtheorem{remark}[theorem]{Remark}
\newcommand{\E}{\mathbb E}
\newcommand{\PP}{\mathbb P}
\newcommand{\R}{\mathbb R}
\newcommand{\norm}[1]{\lVert#1\rVert}
\newcommand{\abs}[1]{\lvert#1\rvert}
\newcommand{\Div}{\operatorname{div}}
\newcommand{\tr}{\operatorname{tr}}
\newcommand{\diag}{\operatorname{diag}}
\newcommand{\supp}{\operatorname{supp}}
\newcommand{\one}{\mathbf 1}
\newcommand{\less}{\lesssim}
\allowdisplaybreaks[1]
\numberwithin{equation}{section}
\setlength{\emergencystretch}{2em}
\title{Talagrand's two-cost discrepancy problem}
\author{Michael Goldman \and Francesco Grotto \and Dario Trevisan}
\date{}
\begin{document}
\maketitle

\begin{aiuse}
This document was generated by an LLM through interaction with the authors,
who do not guarantee its correctness. 
It is made publicly available through the GitHub repository~\cite{repository}, 
with the aim of collecting progress towards the understanding of Talagrand's Ultimate Matching Problem
and its AI-assisted solution. 
It is meant to remain a working note and it elaborates on extensive previous works that are not properly referenced. 
As for all the material in the repository, the authors plan to provide a fully digested
version of the arguments. 
\end{aiuse}

\begin{abstract}
This note addresses Research Problem 17.2.2 in the second edition of
M.~Talagrand's \emph{Upper and Lower Bounds for Stochastic Processes:
Decomposition Theorems}.
We give a necessary and sufficient condition for the discrepancy bound
for every pair of finite convex Young costs. With
$R_i(j)=(\theta_i^*)^{-1}(\exp(2^j))$, the condition is that, for some
$c>0$, both prefix sums
$\sum_{j\le J}\exp[-c R_1(j)^2R_2(J)^2/2^j]$ and their coordinate
interchanges remain bounded uniformly in $J$. 
The proof balances the two coordinate costs across scales through an
adaptive Gaussian construction, with matching lower bounds obtained
from Brownian local time.
A nonmeasurable counterexample explains why completely unrestricted
functions cannot simply be replaced by their convex envelopes.
We thank M. Talagrand for pointing us towards this technical problem as a way to gain better insight on the Ultimate Matching Problem.
\end{abstract}

\section{Statements and the scope of the answer}
Let $n=2^p$ and $G_n=\{1,\ldots,n\}^2$. Differences $D_i$ are taken
only over internal edges. Put
\[
 \mathcal H_n(\theta_1,\theta_2)=
 \left\{h:G_n\to\R:
 \sum\theta_1(|D_1h|)+\sum\theta_2(|D_2h|)\le n^2\right\}.
\]
For independent standard normal variables $g_v$ define
\[
 w_n=\E\sup_{h\in\mathcal H_n}
       \left|\sum_{v\in G_n}g_v(h(v)-\bar h)\right|,
 \qquad \bar h=n^{-2}\sum_vh(v).
\]
The empirical setting is the one in the book: integers
$m_0\le n(v)\le2m_0$, $N=\sum_vn(v)$,
$\mu(v)=n(v)/N$, and independent $U_i$ of law $\mu$.
The relevant density regime is $m_0\ge c_0p$, where $c_0>0$ is fixed.
All constants below are uniform in $p,m_0$, and the admissible counts;
they may depend on the two costs and on $c_0$.
The discrepancy of an empty admissible class is understood as zero.

\begin{theorem}[Two Young costs]\label{pair:thm:criterion}
Suppose $\theta_i:[0,\infty)\to[0,\infty)$ are finite, convex,
$\theta_i(0)=0$, and $\theta_i(t)\ge t$. Let $\psi_i=\theta_i^*$
and let $\rho_i=\psi_i^{-1}$ be the increasing generalized inverses.
Put $R_i(j)=\rho_i(\exp(2^j))$. The following are equivalent:
\begin{enumerate}
\item $w_{2^p}\le C2^{2p}\sqrt p$ for every $p\ge1$.
\item For some $c>0$,
\begin{equation}\label{pair:eq:main-prefix}
 \sup_{J\ge1}\sum_{j=1}^J
 \exp\left[-c\frac{R_1(j)^2R_2(J)^2}{2^j}\right]<\infty,
 \qquad
 \sup_{J\ge1}\sum_{j=1}^J
 \exp\left[-c\frac{R_2(j)^2R_1(J)^2}{2^j}\right]<\infty.
\end{equation}
\item The estimate (17.18) holds uniformly in the above empirical setting:
\[
 \E\sup_{h\in\mathcal H_n}
 \left|\sum_{i=1}^N(h(U_i)-\mu h)\right|
 \le C\sqrt{pm_0}\,n^2.
\]
\end{enumerate}
The same criterion results on replacing $R_i(j)$ by
$e^{2^j}/\theta_i^{-1}(e^{2^j})$. In (2), the two suprema can
equivalently both be required to be at most $e^{-1}$ after increasing $c$.
\end{theorem}

The condition is symmetric. It includes growth profiles whose relative
sizes change infinitely often; no regular variation or common preferred
coordinate is assumed. Its prefix structure is essential. The simpler
pointwise condition $R_1(J)R_2(J)\gtrsim\sqrt{2^J}$ is necessary but
not sufficient.

\begin{theorem}[Finite Borel costs]\label{pair:thm:borel}
Let $\theta_i:[0,\infty)\to[0,\infty)$ be finite Borel functions with
$\theta_i(t)\ge t$. Let $\phi_i$ be the greatest convex minorant on
$\R$ of $s\mapsto\theta_i(|s|)$, and put
\[
 a_i=\inf\theta_i=\phi_i(0),\qquad a=a_1+a_2,
 \qquad \vartheta_i(t)=\phi_i(t)-a_i+t.
\]
If $a\ge1$, the desired estimate always holds. If $a<1$, it holds if
and only if \eqref{pair:eq:main-prefix} holds for the normalized Young
pair $(\vartheta_1,\vartheta_2)$.
\end{theorem}

Every finite nondecreasing cost is Borel and locally bounded, so this
covers arbitrary pairs in the usual monotone-cost interpretation of the
book's question. The Borel assumption can also be replaced by finite
Lebesgue measurability, by completing the finite product measures in the
exceptional-edge argument. A cost that is locally bounded need not be
measurable at all; the same classification still applies, as explained
in the relaxation proof. Completely unrestricted functions require
care: the final section gives an explicit pair for which separate convex
envelopes give the wrong answer. We do not claim a classification of
all such nonmeasurable pairs.

\subsection{Reading the criterion}
If $\theta_1(t)=t$, then $R_1(j)=1$ and the theorem reduces exactly to
\[
 \exists c>0:\qquad
 \sum_{j\ge1}\exp\left[-c\frac{R_2(j)^2}{2^j}\right]<\infty.
\]
Equivalently,
\[
 \exists c>0:\qquad
 \int_{e^e}^{\infty}
 \exp\left[-c\frac{s^2}{\theta_2^{-1}(s)^2\log s}\right]
 \frac{ds}{s\log s}<\infty.
\]
Indeed, the first prefix condition has geometrically decaying tails once
$R_2(J)^2/2^J$ is bounded below; the second becomes the displayed
infinite series. The inverse-product inequality below and dyadic
integration give the integral form.
Consequently the sharp borderline example in this direction is
$t\sqrt{\log t\,\log\log\log t}$; $t\sqrt{\log t}$ fails.
For equal normalized Young costs the exact condition is simply
$\theta(t)\gtrsim t(\log t)^{1/4}$.
Further exact examples are proved after the main argument.

\begin{lemma}[Complementary inverses]\label{pair:lem:inverse-product}
For a finite convex Young function $\theta\ge t$ and its complement
$\psi$, at every $v>0$,
\[
 v\le\theta^{-1}(v)\psi^{-1}(v)\le2v.
\]
\end{lemma}
\begin{proof}
Young's inequality gives the upper bound. For the lower bound set
$x=\theta^{-1}(v)$ and $y=v/x$. If $z\le x$, then
$yz-\theta(z)\le v$. If $z\ge x$, convexity gives
$\theta(z)/z\ge\theta(x)/x=y$. Thus $\psi(y)\le v$ and
$\psi^{-1}(v)\ge y$. Limits give the same assertions for generalized
inverses when the complement has a zero interval or is extended valued.
\end{proof}


% BEGIN elementary_comparisons.tex
\subsection{Conventions and elementary Gaussian comparisons}
On the unit square, $W$ is the isonormal Gaussian process on $L^2$:
$\E W(f)W(g)=\int fg$. Write
$W^0(f)=W(f-\int f)$. If $P_h$ is cell averaging on the dyadic
$h$-square grid, its histogram density has value $h^{-2}W^0(Q)$ on
each cell $Q$. A zero-normal current of divergence $f$ means an
integrable vector field $J$ such that
$\int f\varphi=-\int J\cdot\nabla\varphi$ for smooth tests on the
closed square. For a convex complementary cost $\psi$, use the
Luxemburg norm
\[
 \|g\|_{L^\psi}=\inf\{r>0:\int\psi(|g|/r)\le1\}.
\]
It is a norm, including the extended-valued case: convexity proves
the triangle inequality by taking the convex combination with weights
$r/(r+s)$ and $s/(r+s)$. Young's pointwise inequality gives
$\int|fg|\le2\|f\|_{L^\theta}\|g\|_{L^\psi}$.

\begin{lemma}[Gaussian comparison]\label{pair:lem:gaussian-comparison}
If two finite centered Gaussian processes satisfy
$\E(X_i-X_j)^2\le\E(Y_i-Y_j)^2$ for every $i,j$, then
$\E\max_iX_i\le\E\max_iY_i$. A centered Gaussian covariance
dominated by another has smaller expectation for every convex function
for which the expectations exist. Finally, for deterministic vectors
$a_i$ and a convex positively homogeneous function $F$,
\[
 \E_\varepsilon F\left(\sum_i\varepsilon_i a_i\right)
 \le\sqrt{\pi/2}\,\E_g F\left(\sum_i g_i a_i\right).
\]
\end{lemma}
\begin{proof}
For the first assertion interpolate independent processes by
$Z_t=\sqrt t X+\sqrt{1-t}Y$, and apply Gaussian integration by
parts to $f(z)=\beta^{-1}\log\sum_i e^{\beta z_i}$.
Writing $w_i=e^{\beta z_i}/\sum_j e^{\beta z_j}$, its derivative is
\[
 \frac d{dt}\E f(Z_t)=\frac\beta4\E\sum_{i,j}w_iw_j
 \{\E(X_i-X_j)^2-\E(Y_i-Y_j)^2\}\le0.
\]
Let $\beta\to\infty$. Infinite compact index sets follow by finite
approximations. For covariance domination, add an independent centered
Gaussian of the covariance difference and use conditional Jensen.
For the last assertion write $g_i=\varepsilon_i|g_i|$, and condition
on their independent signs. Jensen over the absolute values gives
the right side at least $\E|g_1|$ times the left side.
\end{proof}

Symmetrization used below also follows directly from Jensen: introduce
independent observations $U_i'$, bound the original centered support
function by its difference-sample expectation, exchange each pair using
an independent sign, and apply the triangle inequality. The result is
twice the signed-process expectation. Subtracting $\mu h$ from every
test before this operation removes the otherwise unbounded constants.

The one-dimensional Gaussian bounds used throughout are
$\PP(|g|\ge t)\le2e^{-t^2/2}$ and
$\PP(|g|\ge t)\ge c e^{-t^2}$ for $t\ge0$. For the second, integrate
the Gaussian density from $t$ to $t+(1+t)^{-1}$ and adjust the fixed
constant on $[0,1]$. Independence then gives
$\E\max_{j\le M}|g_j|\ge c\sqrt{\log(M+1)}$; a union bound gives
the matching upper estimate. Stop-loss comparisons imply convex
order because a convex piecewise linear function is an affine
function plus a nonnegative combination of $(x-t)_+$; approximation
extends this to convex Lipschitz functions. All support functions
used in the finite-grid comparisons are of this type.

% END elementary_comparisons.tex


% BEGIN common_reductions.tex
\section{Gaussian, continuum, and empirical reductions}
\subsection{Grid and continuum comparison}
\label{subsec:grid-continuum-core}

For this subsection assume the two costs are finite convex Young functions, zero at the origin and at least the identity. Write $L=2^p$, $Q=L^2$, and $h=L^{-1}$. On $G_L=\{1,\ldots,L\}^2$
use the modular
\[
 \mathcal E_L(f)=\sum_{\text{horizontal edges}}\theta_1(|\Delta f|)
                    +\sum_{\text{vertical edges}}\theta_2(|\Delta f|).
\]
All sums here contain only internal edges; no boundary edges are added.
For a continuum function $u$ on $[0,1]^2$, let
$f_v=L\,\operatorname{av}_{C_v}u$, where $C_v$ are the $h$ squares.
Jensen along a segment joining adjacent squares gives
\begin{equation}\label{eq:averaging-energy-core}
 \mathcal E_L(f)\le L^2\int_{[0,1]^2}
               \{\theta_1(|u_x|)+\theta_2(|u_y|)\}.
\end{equation}
Indeed each segment average has total overlap at most one after summing
over internal edges. If $W(C_v)=h g_v$ with independent standard
Gaussians $g_v$, then, exactly,
\begin{equation}\label{eq:averaging-pairing-core}
 \sum_v g_v(f_v-\overline f)
       =L^2W(P_h(u-\overline u)).
\end{equation}
Thus a continuum histogram lower bound transfers to the grid without
an approximation or an extra boundary cost.

For the reverse direction, interpolate $h f_v$ bilinearly between cell
centres, and extend constantly in the outward normal direction across
the boundary half strips. Call the result $u_f$. Each derivative is a
convex combination of the corresponding grid differences, giving
\[
 \int(\theta_1(|(u_f)_x|)+\theta_2(|(u_f)_y|))\le h^2\mathcal E_L(f).
\]
The nodal hat functions all have integral $h^2$, including at the
boundary with this extension. Consequently
$\overline{u_f}=h\overline f$, and the elementary mass-matrix bounds give
\[
 c h^4\sum_v(f_v-\overline f)^2
 \le\|u_f-\overline{u_f}\|_2^2
 \le C h^4\sum_v(f_v-\overline f)^2.
\]
Choose a fixed sufficiently fine dyadic subdivision $\delta=h/M$.
The polynomial inverse inequality
$\|\nabla u_f\|_2\le C h^{-1}\|u_f\|_2$, followed by cell Poincare,
shows that $\|P_\delta u_f\|_2\ge c\|u_f\|_2$ for a fixed large $M$.
The same holds for differences of interpolants and after subtracting
constants. Gaussian comparison therefore bounds the grid Gaussian width
by $CL^2$ times the continuum histogram width at mesh $\delta$.
Since $\log(1/\delta)=\log L+O(1)$, this preserves the required scale.

\subsection{Gaussian width and the empirical process}
\label{subsec:empirical-core}

The following comparisons apply to any symmetric, constant-invariant class of lattice functions with bounded diameter modulo constants, including every class considered here. They require no regularity of the original costs.
Assume the law $\mu$ on $G_L$ satisfies
$\max_v\mu(v)\le A/Q$, that $N\max_v\mu(v)\le A m_0$, and
$m_0\ge c_0\log(2Q)$. Constants in this subsection may also depend
on this fixed positive $c_0$. Let $\mathcal H_L=\{f:\mathcal E_L(f)\le Q\}$,
and set
\[
 w_L=\mathbb E\sup_{f\in\mathcal H_L}
            \left|\sum_vg_v(f_v-\overline f)\right|.
\]
Then
\begin{equation}\label{eq:empirical-gaussian-core}
 \mathbb E\sup_{f\in\mathcal H_L}
 \left|\sum_{i=1}^N(f(U_i)-\mu f)\right|
 \le C_A\sqrt{m_0}\,w_L.
\end{equation}
To prove this, first replace each $f$ by $f-\mu f$; constants must be
removed before symmetrization. Conditional on the occupancy counts
$N_v$, symmetrization and the Rademacher-to-Gaussian comparison give a
constant times the Gaussian process with coordinates
$\sqrt{N_v}\,g_v$. Its increment metric is bounded by
$\sqrt{\max_vN_v}\,\|P_\mu(f-f')\|_2$, where
$P_\mu f=f-(\mu f)\mathbf1$. On vectors of uniform mean zero,
\[
 \|P_\mu f\|_2\le(1+\sqrt Q\|\mu\|_2)\|f\|_2
                          \le(1+\sqrt A)\|f\|_2.
\]
Gaussian comparison now bounds its width by
$C_A\sqrt{\max_vN_v}\,w_L$. Finally the exponential moment bound for
a binomial count and a union bound imply
$\mathbb E\max_vN_v\le C_A(m_0+\log(2Q))\le C_{A,c_0}m_0$.
Cauchy--Schwarz proves \eqref{eq:empirical-gaussian-core}.

There is also a converse for the uniform law, valid already for $N\ge Q$:
\begin{equation}\label{eq:empirical-gaussian-lower-core}
 \mathbb E\sup_{f\in\mathcal H_L}
 \left|\sum_{i=1}^N(f(U_i)-\overline f)\right|
 \ge c\sqrt{N/Q}\,w_L.
\end{equation}
Here are details, since this avoids any limit in the sampling density.
For $\lambda\ge1$, let $X=P-P'$ with independent Poisson variables of
mean $\lambda$. There are universal constants $c,C>0$ such that
\[
 \mathbb P(X\ge t)\ge c\exp(-Ct^2/\lambda),\qquad t\ge\sqrt\lambda.
\]
Indeed $\mathbb P(P'\le\lambda)\ge c$. Stirling's formula, summed over
$\lfloor\sqrt\lambda\rfloor$ consecutive integers beginning at
$\lceil\lambda+t\rceil$, gives the claimed lower bound for
$\mathbb P(P\ge\lambda+t)$. The prefactor is bounded below by
$c(1+t/\lambda)^{-1/2}$ and is absorbed in the exponential since
$t\ge\sqrt\lambda\ge1$. Consequently, for $t\ge0$,
\[
 \mathbb E(X-t)_+\ge c\sqrt\lambda
                    \exp[-C(t/\sqrt\lambda+1)^2].
\]
For a sufficiently small universal $a>0$, this dominates
\[
 \mathbb E(a\sqrt\lambda G-t)_+
 \le\frac{a\sqrt\lambda}{\sqrt{2\pi}}
                         e^{-t^2/(2a^2\lambda)}.
\]
Symmetry gives the same stop-loss comparison for negative $t$.
Thus $X$ dominates $a\sqrt\lambda G$ in convex order. Independence
and coordinate-by-coordinate replacement extend this comparison to
any convex function of finitely many coordinates, in particular the
support function of the centered class $\mathcal H_L$.

Put $\lambda=N/Q$ and Poissonize the number of observations to
$M\sim\operatorname{Pois}(N)$. The cell counts $P_v$ are independent
$\operatorname{Pois}(\lambda)$. On uniform-mean-zero tests, the
Poissonized support function is the same for $(P_v)$ and
$(P_v-\lambda)$. Its independent difference has coordinates distributed
as $X$, so the triangle inequality and the convex-order comparison give
\[
 2\,\mathbb E\Big\|\sum_{i=1}^{M}(\delta_{U_i}-\mu)\Big\|_{\mathcal H_L}
 \ge c\sqrt\lambda\,w_L.
\]
If $F(k)$ denotes the expected support function with exactly $k$
observations, centered independent increments and Jensen imply that
$F$ is nondecreasing, while the triangle inequality implies
$F(k+\ell)\le F(k)+F(\ell)$. Hence
$F(k)\le\lceil k/N\rceil F(N)$ and
$\mathbb EF(M)\le2F(N)$. This proves
\eqref{eq:empirical-gaussian-lower-core}. In conjunction with
\eqref{eq:empirical-gaussian-core}, the Gaussian and empirical criteria
are therefore equivalent in Talagrand's regime, without requiring
$m_0\to\infty$ on a fixed grid.

% END common_reductions.tex



% BEGIN selfcontained_local_tools.tex
\section{Local geometric and Gaussian constructions}
\label{sec:selfcontained-local-tools}

This section proves the local facts used by the adaptive flow construction.
The geometry is described in a determinant-one chart. Thus chart area is
physical area. A parent is a parallelogram of diameter $D$ containing a
ball of radius $c_*D$, with fixed $c_*>0$. Terminal pieces are not refined.
All constants below may depend on $c_*$ and on fixed contraction and
H\"older exponents, but not on the mesh or the chart eccentricity in
physical coordinates.

\begin{lemma}[Gaussian concentration]\label{lem:local-gaussian-concentration}
If $F$ is an $L$-Lipschitz function of a standard Gaussian vector, then
\[
 \E e^{t(F-\E F)}\le e^{t^2L^2/2},\qquad
 \PP\{F>\E F+u\}\le e^{-u^2/(2L^2)}.
\]
The same holds for a supremum of Gaussian linear tests whose test vectors
have norm at most $L$, whenever its expectation is finite.
\end{lemma}
\begin{proof}
For a smooth positive $f$ and the Gaussian semigroup
$P_tf(x)=\E f(e^{-t}x+\sqrt{1-e^{-2t}}Z)$, differentiation and Gaussian
integration by parts give
\[
 \operatorname{Ent}_\gamma(f)
 =\int_0^\infty\E\frac{|\nabla P_tf|^2}{P_tf}\,dt.
\]
One obtains the identity first for bounded positive smooth functions by
differentiating $\E[(P_tf)\log(P_tf)]$; its limit at infinity is
$(\E f)\log(\E f)$. Since
$\nabla P_tf=e^{-t}P_t\nabla f$, Cauchy--Schwarz under $P_t$ gives
\[
 \operatorname{Ent}_\gamma(f)
 \le\frac12\E\frac{|\nabla f|^2}{f}.
\]
Apply this to $f=e^{tF}$ and use $|\nabla F|\le L$. If
$A(t)=\log\E e^{tF}$, it follows that
$tA'(t)-A(t)\le t^2L^2/2$. Integrating the derivative of $A(t)/t$,
whose limit at zero is $\E F$, proves the moment bound. Exponential
Markov gives the tail. Smooth approximation and truncation extend the
result to Lipschitz functions. Finite maxima of linear tests are
Lipschitz with the stated constant; increasing finite approximations
and then truncation prove the last assertion.
\end{proof}

\begin{lemma}[A subdivision with a minimum side]\label{lem:local-whitney}
Fix a sufficiently small constant $\tau>0$ and a minimum side $s>0$.
If $\tau D<Cs$, the parent may be declared terminal. Otherwise it has a
measurable subdivision into square children and terminal children. Give a
square its side $l_S$, and a terminal child the label $l_S=s$. Then
\[
 |S|\asymp l_S^2,\quad \operatorname{diam}S\lesssim l_S,
 \quad l_S\asymp\max\{s,\min(\tau D,d_Q(x))\}\quad(x\in S),
 \qquad d_Q(x)=\operatorname{dist}(x,\partial Q).
\]
For a sufficiently small fixed $c>0$, the neighborhoods
$S^+=\{x\in Q:\operatorname{dist}(x,S)<cl_S\}$ have bounded overlap;
intersecting neighborhoods have comparable labels. If $l<l'$ are labels
occurring in one parent, the children of label $l$ lie in a boundary strip
of width $Cl$. A ball of radius $Cl'$ therefore meets at most $C'l'/l$
such children. At a fixed larger label, only $O(1)$ corresponding
neighborhoods can meet a given smaller neighborhood.

The square children contract by an arbitrarily small fixed factor when
$\tau$ is chosen small. After any chart change whose relative metric
eigenvalues lie in a fixed compact interval, the same construction has
uniform constants.
\end{lemma}
\begin{proof}
Use a dyadic lattice of smallest side comparable to $s$; changing $s$ by
a factor two is harmless. Select maximal lattice squares of side
$2s\le l\le\tau D$ contained in $Q$ and satisfying
$\operatorname{dist}(S,\partial Q)\ge\Gamma l$, where $\Gamma$ is a
large fixed constant. Their interiors are disjoint. Every remaining
piece lies in a smallest lattice square, either a full square in $Q$ or
a clipped one. A full uncovered square has distance $O(s)$ from the
boundary: its $2s$ ancestor is under the upper cap and otherwise would
satisfy both containment and clearance.

There is a full uncovered smallest square within $O(s)$ of every clipped
one. To see this, choose a point $x$ in the clipped piece and move a
fraction $As/D$ toward an inball center. Convexity supplies a ball of
radius $Ac_*s$ at the resulting point. Choose $A$ large enough to contain
its lattice square and choose the stopping constant large enough that
the fraction is less than one. The resulting full square is still within
$O(s)$ of the boundary. Taking $\Gamma$ larger than this fixed constant
ensures it is uncovered. Assign each clipped piece to one such full
square, breaking ties in a fixed lattice order. Only $O(1)$ pieces can be
assigned to one receiver, because their lattice squares lie within
$O(s)$ of it. The resulting terminal children have area comparable to
$s^2$ and diameter $O(s)$. They partition the uncovered region.

For a selected square, its dyadic parent fails containment or clearance,
or exceeds the upper cap. In the first two cases $d_Q\lesssim l_S$
throughout the square; in the last $l_S>\tau D/2$. Clearance gives the
converse bound. Together with the terminal estimates this proves the
displayed comparison. The function
$x\mapsto\max\{s,\min(\tau D,d_Q(x))\}$ is $1$-Lipschitz. Choosing
$c$ small in terms of the comparison constants shows that intersecting
neighborhoods have comparable labels. Disjoint child areas and packing
then give bounded overlap.

If a smaller label $l<l'$ occurs, the parent of an $l$ square cannot fail
the upper cap, since $2l\le l'\le\tau D$. It must consequently lie in a
boundary strip of width $Cl$; the same is true of terminal children.
Inside a ball of radius $Cl'$, the four straight boundary strips have
total area $O(ll')$. Dividing by child area $\asymp l^2$ proves the
$O(l'/l)$ count. For a fixed larger label, disjoint areas of size
$\asymp(l')^2$ give the $O(1)$ reverse count. Fixed enlargements only
change the constants. All tests, lattice choices and tie rules are
measurable. A bounded relative chart change preserves an inball and
diameter comparison, proving the final assertion.
\end{proof}

\begin{lemma}[Predictable bridges and a local ray inverse]
\label{lem:local-ray}
Choose each subdivision from the previously revealed Gaussian masses.
Conditionally on this history, the new child-average contrast has its
orthogonal Gaussian law and the remaining noises are independent
mean-zero white-noise bridges in the children.

For such a subdivision there is a continuous zero-normal current $V$
whose divergence is the child-average contrast. On every $S^+$ there is
an even envelope $E_S$ controlling
\[
 \sup_{S^+}|V|+l_S^\zeta[V]_{C^\zeta(S^+)}
             +l_S\|\operatorname{div}V\|_{L^\infty(S^+)},
 \qquad 0<\zeta<1,
\]
with
\[
 \E(E_S\mid\mathrm{old})\le C,\qquad
 \PP(E_S>t\mid\mathrm{old})\le Ce^{-ct^2}.
\]
The pointwise chart covariance of $V$ is at most $CI$. Each of the
linear tests defining the envelope, before subtracting the parent
constant, has a kernel bounded by $C/l_S$ and supported within distance
$Cl_S$ of $S$. The envelope is a fixed sum of suprema of such tests.
\end{lemma}
\begin{proof}
The bridge assertion is orthogonal Gaussian projection. Given the old
history, the child indicators are fixed tests in the old parent bridge.
Their centered span and its orthogonal complement are independent
Gaussian subspaces; the complement is the direct sum of child bridges.
Induction proves the assertion for predictable choices. One may construct
the conditional laws on a countable dense collection of tests and extend
by mean-square continuity; measurability follows from the measurable
indicator maps of the subdivisions.

Let $\eta_Q$ be a smooth probability density in a central ball of $Q$,
with $\|\eta_Q\|_\infty\lesssim D^{-2}$ and
$\|\nabla\eta_Q\|_\infty\lesssim D^{-3}$. Set
\[
 K_Q(x,y)=(x-y)\int_1^\infty t\eta_Q(y+t(x-y))\,dt.
\]
Transport a unit source at $y$ along its straight segment to a point
$z$ averaged with density $\eta_Q(z)$. Its divergence is
$\delta_y-\eta_Q$, and changing variables along the segment gives this
kernel. Convexity keeps the segment in $Q$; extension by zero has no
outer normal flux. Direct differentiation, using
$t|x-y|\le CD$ wherever the integrand is nonzero, gives
\[
 |K_Q(x,y)|\lesssim |x-y|^{-1},\qquad
 |\nabla_xK_Q(x,y)|+|\nabla_yK_Q(x,y)|\lesssim |x-y|^{-2}.
\]
There is also the useful support property
\[
 K_Q(x,y)=0\quad\text{if }|x-y|>Cd_Q(x).
\]
Indeed $x=(1-t^{-1})y+t^{-1}z$, where $z$ is in the central ball. A
nearest supporting line to $x$ has distance at least $cD$ from $z$ and
nonnegative distance from $y$, so $d_Q(x)\ge cD/t$, while
$|x-y|\le CD/t$. In particular $\int|K_Q(x,y)|dy\lesssim d_Q(x)$,
which shows that the inverse of a bounded source is continuous and
vanishes on the boundary. More explicitly, splitting into disks of
radius $2a$ about two poles and the remaining annulus yields
\[
 \int|K_Q(x,y)-K_Q(x',y)|dy
 \lesssim a(1+\log(R/a)),\quad a=|x-x'|\le R/2,
\]
when both supports lie in a ball of radius $CR$.

Let $P$ be child averaging and $P_Q$ parent averaging. The contrast
current is
$V(x)=W((P-P_Q)K_Q(x,\cdot))$. Its divergence is the required contrast.
For $x,x'\in S^+$, take $R=l_S$. A child meeting a sufficiently small
ball about $x$ has side comparable to $l_S$ by Lemma~\ref{lem:local-whitney};
averaging the $|x-y|^{-1}$ singularity there gives $C/l_S$. Farther
children satisfy the same bound directly. Every child meeting the kernel
support has diameter $O(l_S)$, so averaging enlarges that support by
only $O(l_S)$. The two-pole estimate, divided by the near child area,
and the derivative bound on far children give, for any $\beta<1$,
\[
 \|PK_Q(x,\cdot)\|_\infty\le C/l_S,
 \quad
 \left\|\left(\frac{l_S}{|x-x'|}\right)^\beta
 P(K_Q(x,\cdot)-K_Q(x',\cdot))\right\|_\infty\le C_\beta/l_S.
\]
These kernels have support area $O(l_S^2)$ and therefore bounded
$L^2$ norms. The normalized density test is
$l_SP(x,\cdot)$; it has exactly the same bounds. Subtracting $P_Q$ is
a common orthogonal projection and contracts $L^2$ norms. This proves
the covariance assertion.

Choose $\zeta<\beta<1$. On rescaling $S^+$ to diameter $O(1)$, nets of
mesh $2^{-j}$ have cardinality $O(2^{2j})$. The maximum of the Gaussian
differences between successive nets and neighboring net points has
expectation $O(2^{-\beta j}\sqrt{1+j})$. Summing after multiplication
by $2^{\zeta j}$ gives finite expected supremum and $C^\zeta$ seminorm.
Only $O(1)$ comparable children meet $S^+$, so its normalized density
maximum has bounded expectation as well. Each envelope test has bounded
$L^2$ norm; Lemma~\ref{lem:local-gaussian-concentration} gives the tail.
Taking absolute suprema makes the envelope even in all fresh Gaussian
coordinates. The localized preprojection kernel bounds just proved apply
to every test in these suprema.
\end{proof}

\begin{lemma}[Joint envelope tails]\label{lem:local-joint-tail}
Conditionally on the old history, for every predictable finite set
$\mathcal A$ of attempted children and all sufficiently large $t$,
\[
 \PP\{E_S>t\text{ for every }S\in\mathcal A\mid\mathrm{old}\}
 \le e^{-c|\mathcal A|t^2}.
\]
\end{lemma}
\begin{proof}
Choose one signed test kernel $k_S$ per child, before subtracting parent
constants. If their supports overlap and $l=l_S\le l_T=l'$, then
\[
 |\langle k_S,k_T\rangle|\le C l/l'.
\]
For a fixed larger child, Lemma~\ref{lem:local-whitney} bounds the number
of overlapping smaller tests at side $l$ by $Cl'/l$. For a fixed
smaller child, only $O(1)$ tests at each larger side can overlap. Apply
the weighted Schur test with weights $l_S^\eta$, $0<\eta<1$. For a row
at the larger scale, the contribution from scale $l$ is at most
$C(l/l')^\eta$; for a row at the smaller scale, the contribution from
$l'$ is at most $C(l/l')^{1-\eta}$. Both geometric sums are bounded.
Comparable scales are controlled by packing. Thus
\[
 \Big\|\sum_Sc_Sk_S\Big\|_2^2\le C\sum_S|c_S|^2
\]
uniformly over all choices of tests. Within each parent, subtracting
its constant is a common orthogonal projection and cannot increase this
norm. Different parents are disjoint Gaussian source blocks, so the
same inequality holds across all parents and their different charts.

The sum of envelopes is a fixed sum of suprema over choices of one
test per child. The preceding bound makes it $C\sqrt{|\mathcal A|}$-
Lipschitz in the fresh Gaussian vector. Its conditional mean is
$O(|\mathcal A|)$. The event in the lemma forces this sum above
$t|\mathcal A|$. Gaussian concentration proves the claim.
\end{proof}

\begin{lemma}[Even clipping and exact omitted sources]
\label{lem:local-clipping}
Fix a threshold $E_0$. The fresh current can be multiplied pointwise by
an even scalar in $[0,1]$ so that every retained childwise current and
scaled H\"older envelope is at most $CE_0$. Its omitted divergence is a
sum of bounded mean-zero sources $f_S$, one per exceptional child, with
\[
 \|f_S\|_\infty\le C\sigma E_S/l_S,
 \qquad \|f_S\|_1\le C\sigma E_Sl_S.
\]
Here $\sigma$ is the scale multiplying the fresh current, and $f_S$ is
supported in the physical image of an intrinsic ball of radius $Cl_S$
intersected with its parent. Its associated omitted current has zero outer
normal trace, and in particular the source has integral zero.
At a fixed point this attenuation does not increase the exponential
moment of any fixed linear functional of the fresh Gaussian current.
\end{lemma}
\begin{proof}
Call $S$ exceptional when $E_S>E_0$. Take a Lipschitz cutoff $\chi_S$
equal to one on $S$, supported in $S^+$, with chart gradient at most
$C/l_S$. Multiply $V$ by
$a=\prod_{S\text{ exceptional}}(1-\chi_S)$.
The factor is even and in $[0,1]$. It vanishes on exceptional children.
On a nonexceptional child, only boundedly many cutoffs occur, at
comparable sides. The product difference inequality and their Lipschitz
bounds show that the retained supremum and scaled H\"older seminorm
are at most $CE_S\le CE_0$.

Order the exceptional children and write
$1-a=\sum_S\varphi_S$, where
$\varphi_S=\chi_S\prod_{T<S}(1-\chi_T)$. Bounded comparable overlap gives
$|\nabla\varphi_S|\le C/l_S$. Define
$f_S=\sigma\Div(\varphi_SV)$. The product rule and the envelope's
density bound give the asserted $L^\infty$ bound; the support has area
$O(l_S^2)$, giving the $L^1$ bound. The uncut ray current is continuous
and vanishes on the parent boundary. Therefore the product's zero
extension has no boundary measure in its divergence, and its integral
is zero. Determinant-one pushforward preserves these statements.

For the moment assertion let $g$ be the entire fresh Gaussian vector
and $\ell(g)$ the fixed linear functional at the point. Since
$a(-g)=a(g)$,
\[
 \E e^{a(g)\ell(g)}=\E\cosh(a(g)\ell(g))
 \le\E\cosh(\ell(g))=\E e^{\ell(g)}.
\]
\end{proof}

\begin{lemma}[Sparse exact repair]\label{lem:local-sparse-repair}
Let $m=\log(1/h)$. Suppose there are at most $Cm$ generations, every
chart and its inverse have norm at most $m^C$, and intrinsic child labels
lie in $[h,m^C]$. The domain is the unit square; the same proof applies
to an area-one axis rectangle with polynomially bounded sides and
reciprocals. Assume
\[
 \log m=o(E_0^2),\qquad E_0^2\le m,
 \quad D_0=m/E_0^2,\qquad T_0=\sigma\sqrt m.
\]
For every $1\le B\le m^C$, with probability $1-o(1)$ all omitted
sources in Lemma~\ref{lem:local-clipping} have a combined exact
zero-normal repair supported on a set of area at most $e^{-cE_0^2}$,
and
\[
 \|J_y\|_\infty\le CT_0D_0/B,
 \qquad \|J_x\|_{\psi_2}\le CT_0D_0B/E_0,
 \qquad \psi_2(t)=e^{t^2}-1.
\]
The conclusion remains valid before any stopping rule based on the
previously revealed history.
\end{lemma}
\begin{proof}
Each child has area at least $ch^2$, so there are at most $Cm h^{-2}$
attempts, up to polynomial factors. The conditional Gaussian tails and
a union bound show $E_S\le C\sqrt m$ for every attempt, with probability
$1-o(1)$. On this event an omitted source in a determinant-one ellipse
of intrinsic radius $Cl$ satisfies
\[
 \|f_S\|_\infty\le CT_0/l,
 \qquad \|f_S\|_1\le CT_0l.
\]
Write $G=AA^{\mathsf T}$ for its chart metric. The maximal horizontal
and vertical chords of this ellipse are bounded by
$Cl/\sqrt{G_{yy}}$ and $Cl/\sqrt{G_{xx}}$. This follows by setting the
other coordinate equal to zero in its quadratic form, since $\det G=1$.
Thus the maximal absolute horizontal and vertical line masses, $A_x,A_y$,
satisfy $A_xA_y\le CT_0^2$.

Choose, before the fresh reveal, an axis box in the domain containing
the support ellipse portion, with both sides at least $Bl$. Enclose
its coordinate projections, enlarge, and slide the intervals inside
the domain. Each side is at most a polynomial times $l$. If such a
box cannot fit, then $l\ge m^{-C'}$; area packing leaves only
polynomially many such candidates per generation. The probability any
of them is exceptional is $o(1)$, since $E_0^2\gg\log m$.

In a box $I\times J$ and for any bounded mean-zero $f$, set
\[
 r(y)=\int_I f(x,y)dx,\qquad
 J_x(x,y)=\int_{\inf I}^x\left(f(t,y)-\frac{r(y)}{|I|}\right)dt,
 \quad J_y(x,y)=\frac1{|I|}\int_{\inf J}^y r(t)dt.
\]
These primitives have divergence $f$ and zero normal boundary trace.
If $A_x\le T_0B$, this gives horizontal supremum $2A_x\le CT_0B$ and
vertical supremum at most $\|f\|_1/|I|\le CT_0/B$. If $A_x>T_0B$,
then $A_y\le CT_0/B$ and the column version gives the same two bounds.
Extending each box current by zero preserves its divergence.

We now control overlap, including dependence across generations. Fix a
point and pad all predictable boxes by $O(h^2)$. In one generation at
most $m^{C'}$ candidate boxes contain the point: at each intrinsic
dyadic size their disjoint child areas lie within a polynomial times
that size of the point, and there are $O(m)$ size classes. Let $I_S$
indicate exceptionality among these candidates. For
$\lambda=cE_0^2/2$, the joint tail lemma and expansion over subsets give
\[
 \E\left(e^{\lambda\sum_S I_S}\mid\mathrm{old}\right)
 =\E\left(\prod_S[1+(e^\lambda-1)I_S]\mid\mathrm{old}\right)
 \le\exp\{m^{C'}e^{-cE_0^2/2}\}.
\]
Iterate over the $O(m)$ generations, adding empty ones after a stopping
rule. Chernoff's inequality bounds the probability of overlap exceeding
$C'D_0$ by $\exp(-c'C'm+o(1))$. A union bound over a mesh of spacing
$h^2$ and the padding gives a uniform overlap bound $CD_0$ for a
sufficiently large fixed $C'$. Padding is harmless in the candidate
count because physical sizes are at least $h$ times an inverse
polynomial, which exceeds $h^2$.

The sum of candidate box areas in one generation is polynomially
bounded: each is at most a polynomial times $l_S^2$, while the child
areas are disjoint. Conditional single-child tails therefore bound the
expected total exceptional-box area over all generations by
$m^{C''}e^{-cE_0^2}$. Markov's inequality makes this area at most
$e^{-c'E_0^2}$ with probability $1-o(1)$. The uniform overlap gives
hard supremum $CT_0D_0/B$ and other supremum $M=CT_0D_0B$ on this
sparse support. A function bounded by $M$ on area $q$ has
\[
 \|f\|_{\psi_2}\le\frac{M}{\sqrt{\log(1+1/q)}}.
\]
Substitution in its Luxemburg modular proves this inequality and the
claimed soft bound. Every repair is exact and zero-normal, so their
sum has the stated divergence.
\end{proof}

\begin{corollary}[Sparse repair in two Orlicz norms]
\label{lem:local-sparse-orlicz}
In the setting of Lemma~\ref{lem:local-sparse-repair}, let $\psi_i$ be
Young functions with generalized inverses $\rho_i$. Suppose
\[
 \rho_1(e^u)\rho_2(e^u)\ge c\sqrt u
 \quad\text{for all sufficiently large }u.
\]
Then, for $E_0=m^{1/2-\eta}$ and $0<\eta<1/10$, the omitted normalized
sources, with $T_0=O(1)$, have an exact repair satisfying
\[
 \max_i\|J_i\|_{L^{\psi_i}}
 \le C m^{(5/2)\eta-1/4}=o(1)
\]
with probability $1-o(1)$.
\end{corollary}
\begin{proof}
The proof of the sparse-repair lemma gives, on area at most
$q=e^{-c_0E_0^2}$, the two supremum bounds
$CT_0D_0B$ and $CT_0D_0/B$. It is symmetric under exchanging axes,
so the same statement holds for polynomially bounded positive $B$ and
$B^{-1}$. Put
\[
 R_i=\min\{\rho_i(e^{c_0E_0^2}),m\},\qquad
 B=\sqrt{R_1/R_2}.
\]
Each $R_i$ is bounded below by a fixed positive constant for large $m$.
If neither truncation is active their product is at least $c'E_0$.
If a truncation is active, their product is at least $c'm\ge c'E_0$.
Thus $B,B^{-1}$ are polynomially bounded and $R_1R_2\ge c'E_0$.
A function bounded by $M$ on area $q$ has Luxemburg norm at most
$M/\rho_i(1/q)$. Consequently both repair norms are at most
\[
 \frac{CT_0D_0}{\sqrt{R_1R_2}}
 \le\frac{Cm}{E_0^{5/2}}
 =Cm^{(5/2)\eta-1/4}.
\]
\end{proof}

\begin{lemma}[Grid capacity and Gaussian repair]\label{lem:local-grid-capacity}
On the graph of a square of $K^2$ cells, with $K$ a power of two, orient internal faces and let
$dq$ be outgoing minus incoming load at each cell. For $\sum_vg_v=0$,
\[
 \mathcal B(g):=\min_{dq=g}\max_e|q_e|
 =\max_{\varnothing\ne S\ne D}
      \frac{|\sum_{v\in S}g_v|}{|\partial S|}
 \le C\|g\|_2.
\]
Consequently $\mathcal B$ is $C$-Lipschitz. If $g$ is centered Gaussian
with covariance at most the projection onto mean-zero vectors, then
\[
 \E\mathcal B(g)\le C(1+\log K),\qquad
 \PP\{\mathcal B(g)>C(1+\log K)+t\}\le e^{-ct^2}.
\]
For physical cell width $w$, interpolating a component linearly between
its two normal face loads gives a zero-normal continuum current with
component suprema at most $\mathcal B(g)$ and divergence $g_v/w$ on cell
$v$.
\end{lemma}
\begin{proof}
Any feasible flow gives the cut lower bound by summing divergence over
$S$. Conversely, the image under $d$ of the face cube $[-b,b]^E$ has
support function $b\sum_{v\sim z}|u_v-u_z|$ at a vertex function $u$.
Subtract $\min u$ and integrate superlevel sets. The cut inequalities
imply $\langle g,u\rangle\le b\sum_{v\sim z}|u_v-u_z|$ for every $u$;
finite-dimensional convex separation therefore puts $g$ in that image.
This proves equality.

For $|S|\le K^2/2$, if no occupied row or column is full, each occupied
row and column is mixed. Their numbers have product at least $|S|$,
so their cut edges give $|\partial S|\ge c\sqrt{|S|}$. If a row is full,
all columns are occupied and at least $K/2$ are mixed, since at most
$|S|/K\le K/2$ columns can be full. The same argument handles a full
column. Thus in all cases
$|\partial S|\ge c\sqrt{\min(|S|,|D\setminus S|)}$.
Cauchy--Schwarz in the cut formula proves the Euclidean bound, and
subadditivity gives its Lipschitz assertion.

For independent standard Gaussians $Z_v$, decompose their centered
vector into orthogonal dyadic child-average contrasts. In a parent of
side $L$ cells, each of its four child averages has standard deviation
$O(1/L)$. Subtract row means and take horizontal primitives, then take
vertical primitives of those row means. The result routes the parent
contrast with zero outer load. Each face load in this one-level flow
has variance $O(1)$, because its coefficients on child values are
$O(L)$. Different levels are orthogonal Gaussian subspaces, so the
summed flow has face variance $O(1+\log K)$. There are $O(K^2)$ faces;
the Gaussian maximum estimate gives expected supremum $C(1+\log K)$.
For covariance dominated by the mean-zero projection, add an independent
Gaussian with the covariance difference and use convexity of
$\mathcal B$ to obtain the same expected bound. Gaussian concentration
and its Euclidean Lipschitz constant prove the tail. Linear normal-face
interpolation has matching traces across internal faces, zero exterior
normal trace, and divergence equal to the face-load difference divided
by $w$, proving the last statement.
\end{proof}

\begin{lemma}[Collected local tools]\label{lem:local-tools}
The ordinary adaptive refinement may use a subdivision with any fixed
sufficiently small intrinsic contraction, a pointwise chart covariance
bound, even child envelopes with uniform Gaussian and joint tails, even
clipping with exact omitted sources, the sparse-repair bounds of
Lemma~\ref{lem:local-sparse-repair}, and the dimension-independent grid
capacity and Gaussian repair bounds of
Lemma~\ref{lem:local-grid-capacity}. All remain valid conditionally on
previous reveals and before predictable stopping rules.
\end{lemma}
\begin{proof}
This is precisely the content of Lemmas~\ref{lem:local-whitney}--
\ref{lem:local-grid-capacity}. Their proofs use only conditional
orthogonal projection, the parent geometry, and deterministic kernel
and packing bounds, so their constants do not depend on the realized
past or the physical eccentricities allowed in the sparse-repair lemma.
\end{proof}

% END selfcontained_local_tools.tex


\section{Constructing the two-cost phase}

% BEGIN pair_width_regularization.tex
\subsection{Regularizing arbitrary Young widths}
\label{pair:sec:width-regularization}

No regular growth assumption on the Young functions is needed. The
following reduction supplies it while shrinking the permissible widths.
For squared widths $b_{i,j}=\rho_i(e^{2^j})^2$ put
\[
 v_{i,j}=\frac{2^{j/2}}{\sqrt{b_{i,j}}},\qquad
 G_{i,J}=\mathbb E\max_{j\le J}v_{i,j}|g_j|.
\]
The two prefix entropy conditions are equivalent, up to their fixed
constants, to
\begin{equation}\label{pair:eq:gaussian-prefix-regularization}
 G_{1,J}\le C\sqrt{b_{2,J}},\qquad
 G_{2,J}\le C\sqrt{b_{1,J}}\qquad(J\ge1).
\end{equation}
First the full family of prefix estimates implies $\inf_jp_j>0$, as
shown in the first paragraph of the allocation lemma below. Thus
$v_{i,j}/\sqrt{b_{3-i,J}}\le C$ for $j\le J$, by monotonicity.
Under this largest-coordinate bound, the standard Gaussian-maximum
equivalence applies: if
$\sum_j\exp(-cR^2/v_j^2)\le C$ and $\max_jv_j\le C_0R$, a union bound
integrated above a sufficiently large multiple of $R$ gives
$\mathbb E\max_jv_j|g_j|\le C'R$. Conversely, if this expectation is
at most $CR$, the indices with $v_j\ge t$ have cardinality at most
$\exp(C'R^2/t^2)$, because the expected maximum of $M$ independent
standard Gaussian absolute values is at least $c\sqrt{\log(M+1)}$.
Partitioning the $v_j$ into dyadic bands gives the exponential sum bound
with a sufficiently large exponent constant. These arguments apply
uniformly to every finite prefix. The harmless largest-coordinate issue
can alternatively be read directly from
$G_{i,J}\ge\sqrt{2/\pi}\,v_{i,j}$.

Fix $q>1$. Regularize one squared width by
\begin{equation}\label{pair:eq:past-width-envelope}
 \widetilde b_{1,j}=\inf_{k\le j}2^{q(j-k)}b_{1,k}.
\end{equation}
It is nondecreasing, satisfies
$1\le\widetilde b_{1,j}\le b_{1,j}$, and
$\widetilde b_{1,j+1}\le2^q\widetilde b_{1,j}$.
We claim that both inequalities in
\eqref{pair:eq:gaussian-prefix-regularization} remain valid with this
new width. Write $r=2^{-(q-1)/2}<1$. Its new reciprocal width is
\[
 \widetilde v_{1,j}=\sup_{k\le j}r^{j-k}v_{1,k}.
\]
Consequently
\begin{align*}
 \mathbb E\max_{j\le J}\widetilde v_{1,j}|g_j|
 &\le\sum_{\ell\ge0}r^\ell
       \mathbb E\max_{k\le J-\ell}v_{1,k}|g_{k+\ell}|\\
 &\le(1-r)^{-1}G_{1,J}\le C\sqrt{b_{2,J}}.
\end{align*}
For the other inequality, take an index $k\le J$ achieving
$\widetilde b_{1,J}=2^{q(J-k)}b_{1,k}$. The old prefix inequalities
imply the pointwise product lower bound
$b_{1,k}b_{2,k}\ge c2^k$. For $\ell>k$, monotonicity of $b_2$ gives
\[
 v_{2,\ell}\le2^{(\ell-k)/2}v_{2,k}
                  \le C2^{(\ell-k)/2}\sqrt{b_{1,k}}.
\]
The expected maximum over these later indices is therefore at most
$C2^{(J-k)/2}\sqrt{b_{1,k}}$, using a geometric sum. The earlier
indices contribute $G_{2,k}\le C\sqrt{b_{1,k}}$. Since $q>1$,
\[
 G_{2,J}\le C2^{q(J-k)/2}\sqrt{b_{1,k}}
                         =C\sqrt{\widetilde b_{1,J}}.
\]
This proves the claim. Applying the same operation to the second width
preserves both inequalities once again.

Interpolate the logarithms of the resulting widths linearly at
$s_j=j\log2$. They are nondecreasing and have Lipschitz constant at most
$q$. If desired, convolve these logarithms with one fixed smooth positive
kernel. This changes each width by at most a fixed factor, preserving
the prefix conditions. At every intermediate $s$, the resulting width
is bounded by a fixed multiple of the original
$\rho_i(e^{e^s})^2$: use the left endpoint of the containing interval,
monotonicity of the original width, and the factor $2^q$ between the two
new endpoints. Thus the regularized widths may be used for an upper
bound for the original Young functions.

% END pair_width_regularization.tex


% BEGIN pair_entropy_path.tex
\subsection{From prefix entropy to smooth squared semiaxes}
\label{pair:sec:entropy-path}

The following elementary construction is useful because it permits the
preferred coordinate direction to change arbitrarily often.
Write $s_j=j\log2$ and
\[
 b_{i,j}=\rho_i(e^{2^j})^2,\qquad
 p_j=b_{1,j}b_{2,j}2^{-j}.
\]
Assume the two prefix estimates
\begin{equation}\label{pair:eq:prefix-for-path}
 \sup_J\sum_{j\le J}
 \exp[-c\,b_{1,j}b_{2,J}2^{-j}]<\infty,
 \qquad
 \sup_J\sum_{j\le J}
 \exp[-c\,b_{2,j}b_{1,J}2^{-j}]<\infty
\end{equation}
for some $c>0$. Increasing $c$ does not invalidate these estimates.
Suppose also that the widths have been regularized in logarithmic scale,
so that there are positive nondecreasing functions $b_i(s)$, dominated by the original Young widths up to fixed factors, satisfying
the same prefix estimates at $s_j$, and for which
$\log b_i$ are Lipschitz. All conclusions below allow fixed factors in
the upper bounds on the widths.

\begin{lemma}[Smooth allocation of the two widths]
\label{pair:lem:smooth-allocation}
For every sufficiently small $\eta>0$ there are smooth positive
functions $a(s),b(s),T(s)$ on a half line $s\ge s_0$, and constants
$c_*,\varepsilon_*,C_*,K>0$, such that
\begin{gather}
 a(s)\le C_*b_1(s),\qquad b(s)\le C_*b_2(s),\qquad
 a(s)b(s)=e^sT(s),                                             \label{pair:eq:allocated-product}\\
 K\le T(s)\le C_*\log(e+s-s_0),\qquad
 |T^{(q)}(s)|\le C_q\eta\quad(q\ge1),                         \label{pair:eq:allocated-clock}\\
 \alpha=(\log a)',\quad\beta=(\log b)',\qquad
 \alpha,\beta\ge\varepsilon_*e^{-c_*T},\qquad
 \alpha+\beta=1+T'/T,                                        \label{pair:eq:allocated-slopes}\\
 |\alpha^{(q)}|\le C_q\alpha,
 \qquad |\beta^{(q)}|\le C_q\beta\quad(q\ge1).                 \label{pair:eq:allocated-relative}
\end{gather}
Here $K$ can be prescribed beforehand. The constants may depend on the
prefix bounds, the width regularization, and the chosen parameters.
The coefficient of $\eta$ in the derivative bounds is fixed; by starting
with a smaller parameter it can be made as small as subsequently needed.
\end{lemma}

\begin{proof}
First, $\inf_jp_j>0$. Indeed, monotonicity of the widths gives, for
$j-r\le k\le j$,
\[
 b_{1,k}b_{2,j}2^{-k}\le 2^r p_j.
\]
If $p_j$ tends to zero, the number of preceding indices for which this
quantity is at most one tends to infinity, contradicting the first
prefix estimate. There are enough preceding indices because the widths
are bounded below by a positive constant. Hence the continuous function
$p(s)=b_1(s)b_2(s)e^{-s}$ is bounded below as well, by local comparability.

For each $i=1,2$, the prefix estimates imply the interval bound
\begin{equation}\label{pair:eq:unweighted-clock}
 \sum_{r\le j\le J}e^{-c_0p_j}
 \le C\left[1+\log\frac{b_{i,J}}{b_{i,r}}\right].
\end{equation}
To prove this, partition the indices by the dyadic bands of $b_{i,j}$.
Within any one band let $J'$ be its last index. The appropriate prefix
estimate at $J'$ bounds the sum over this band of $e^{-2cp_j}$, since
$b_{i,J'}/b_{i,j}\le2$. The number of bands is at most a constant times
$1+\log(b_{i,J}/b_{i,r})$. This proves
\eqref{pair:eq:unweighted-clock} with $c_0=2c$.
Equivalently, the same argument gives the two weighted Carleson bounds
\[
 \sum_{j\le J}\frac{b_{i,j}}{b_{i,J}}e^{-c_0p_j}\le C.
\]

The constants can be chosen without a circular dependence on the desired
small slope $\eta$. First choose the fixed exponential clock coefficient.
Then take $K=K(\eta)$ so large that the factors
$e^{-c_1K}/\eta^r$ occurring below, for the finitely many required powers
$r$, are bounded by fixed constants. The exponential-kernel boundary
sums and the additional caps then have uniform interval-clock constants.
In particular the small multiplier $\varepsilon$ below can be chosen
independently of the subsequently needed smallness of $\eta$. Increasing
$A$ ensures the prescribed lower bound $K$, and only changes fixed width
and Hall padding constants; such constants are permitted here.

Choose a large constant $A$ and define the Lipschitz lower envelope
\[
 T_0(s)=\inf_{j\ge1}\{A p_j+\eta|s-s_j|\}.
\]
By increasing $A$, its lower bound is at least the prescribed $K$;
local comparability gives $T_0(s)\le C p(s)$. Its Lipschitz constant is
at most $\eta$. Choose constants $A_1,R$ with $A_1c_1>2$ and
$R\ge A_1/\eta$, where $c_1>0$ is fixed, and put
\[
 T_\mathrm{cap}(s)=K+A_1\log(1+(s-s_0)/R),\qquad
 T=\min\{T_0,T_\mathrm{cap}\}.
\]
Adjusting $s_0$ by a fixed amount, if necessary, leaves all estimates
unchanged. This cap is also $\eta$-Lipschitz and
$\int_{s_0}^\infty e^{-c_1T_\mathrm{cap}}<\infty$.
Choose $A$ so that $c_1A\ge c_0$. Since
\[
 e^{-c_1T_0(s)}
 \le\sum_j e^{-c_1A p_j}e^{-c_1\eta|s-s_j|},
\]
integration over an interval, followed by
\eqref{pair:eq:unweighted-clock}, yields for both $i$
\begin{equation}\label{pair:eq:continuous-clock}
 \int_v^w e^{-c_1T(s)}\,ds
 \le C\left[1+\log\frac{b_i(w)}{b_i(v)}\right],\qquad w\ge v.
\end{equation}
The indices outside the interval contribute only a fixed geometric
boundary sum. The logarithmic Lipschitz bound on the widths absorbs the
fixed enlargement of its endpoints. The cap contributes a finite global
constant.

Set $m(s)=\varepsilon e^{-c_1T(s)}$, with $\varepsilon>0$ sufficiently
small, and $\tau(s)=s+\log T(s)$. Then
$\tau'\ge9/10$, $2m<\tau'$, and, for both $i$,
\begin{equation}\label{pair:eq:small-clock}
 \int_v^w m\le C_0+\tfrac12\log\frac{b_i(w)}{b_i(v)}.
\end{equation}
Because $T(v)\le C p(v)$, a fixed $D>0$ consequently ensures the two
Hall inequalities
\begin{equation}\label{pair:eq:hall-allocation}
 \tau(v)+\int_v^w m
 \le \log b_1(v)+\log b_2(w)+2D,
 \qquad
 \tau(v)+\int_v^w m
 \le \log b_2(v)+\log b_1(w)+2D
\end{equation}
for every $v\le w$.

Here is an explicit bounded-slope selection, avoiding a compactness
argument. Define
\[
 M(s)=\int_{s_0}^s m,
 \quad N(s)=\tau(s)-\tau(s_0)-2M(s),
 \quad L(s)=\tau(s)-\log b_2(s)-D,
 \quad U(s)=\log b_1(s)+D.
\]
The function $N$ is strictly increasing. Put
\[
 A(s)=M(s)+\sup_{t\ge s_0}
 \{L(t)-M(t)-(N(t)-N(s))_+\}.
\]
The supremum is finite and $L\le A\le U$: the required pairwise
inequalities are exactly \eqref{pair:eq:hall-allocation}. As a function
of $N(s)$, the supremum is nondecreasing and $1$-Lipschitz.
Consequently $m\le A'\le\tau'-m$ almost everywhere. Put
$B=\tau-A$, $a=e^A$, $b=e^B$. Then
$a\le e^D b_1$, $b\le e^D b_2$, $ab=e^sT$, and both logarithmic
slopes are at least $m$.

It remains to smooth without losing these useful properties. Extend
$A,B$ to $s<s_0$ with slopes $1/2$ and extend $T$ constantly. Convolve
$A$ and $B$ with the same even probability kernel
\[
 k(s)=Z^{-1}e^{-\sqrt{1+s^2}},
 \qquad |k^{(q)}|\le C_q k.
\]
Write $\widetilde A=k*A$, $\widetilde B=k*B$ and
$\widetilde T=\exp[k*\log T]$. Since $k$ has mean zero,
$\widetilde A+\widetilde B=s+\log\widetilde T$.
The original logarithmic slopes are nonnegative and bounded. Hence
convolution changes $A,B$ by bounded constants, preserves their width
upper bounds up to fixed factors, and gives
\[
 |(\widetilde A')^{(q)}|\le C_q\widetilde A',\qquad
 |(\widetilde B')^{(q)}|\le C_q\widetilde B'.
\]
The Lipschitz property of $T$ implies
$|\widetilde T-T|\le C\eta$ and
$|\widetilde T^{(q)}|\le C_q\eta$.
For completeness, these estimates follow by splitting the exponentially
weighted convolution at $|t|=T(s)/(2\eta)$; on the inner part use
$T(s-t)=T(s)+O(\eta|t|)$, and on the outer part use $T\ge K$ and the
exponential kernel tail. Derivatives of $\log\widetilde T$ are obtained
by convolving $(\log T)'$ with derivatives of $k$.
Also
\[
 \widetilde A'(s)\ge
 \varepsilon e^{-c_1T(s)}\int k(t)e^{-c_1\eta|t|}\,dt
 \ge\varepsilon_*e^{-c_*\widetilde T(s)},
\]
and similarly for $\widetilde B'$. The logarithmic upper bound on $T$
is preserved by convolution. Dropping tildes proves the lemma.

For the quantitative spatial construction one may, uniformly in a mesh
parameter $m$, include the additional cap
\[
 T_{\mathrm{end},m}(s)=K+\eta\sqrt{1+(s-\log m)^2}
\]
when taking the minimum defining $T$. Its exponential clock has a
uniformly bounded integral, so every Hall estimate and every constant
above is uniform in $m$. The smoothing changes this upper bound by at
most a fixed additive constant. In particular, for $1\le u\le C_1m$,
\begin{equation}\label{pair:eq:mesh-adapted-clock}
 e^{C T_m(\log u)}\le C'\,(m/u)^{C'\eta},
\end{equation}
where the constants are independent of $m$. At the left endpoint,
$T_m(s_0)\le T_0(s_0)$; the two width bounds and their product therefore
also give fixed positive lower bounds on the initial squared semiaxes.
\end{proof}

% END pair_entropy_path.tex


% BEGIN rounded_rectangle_phase.tex
\subsection{A rounded-rectangle Bellman supersolution}
\label{pair:sec:rounded-phase}

We use the functions supplied by Lemma~\ref{pair:lem:smooth-allocation},
with $s=\log u$. Thus $a,b$ are squared semiaxes and
\[
 ab=uT,\quad \alpha=D\log a,\quad\beta=D\log b,\quad
 \gamma=D\log T,\quad\alpha+\beta=1+\gamma,
 \qquad D=u\frac d{du}.
\]
They satisfy, uniformly in the mesh parameter when it is present,
\begin{gather*}
 T\ge K,\qquad |D^qT|\le C_q\eta,\qquad
 \alpha,\beta\ge\delta:=\varepsilon_*e^{-c_*T},\\
 |D^q\alpha|\le C_q\alpha,\quad
 |D^q\beta|\le C_q\beta,
 \qquad T\le C\log(e+\log u).
\end{gather*}
The small parameter $\eta$ is chosen after the fixed constants below;
the lower bound $K$ can be increased as described in the allocation
construction. Choose $\Lambda>c_*+1$ and set
\[
 P=\Lambda T,\quad \lambda=P/a,\quad\mu=P/b,
 \quad X=\lambda x^2,\quad Y=\mu y^2,
 \quad E=e^X,\quad F=e^Y.
\]
For each $u$, consider the convex oval
\begin{equation}\label{pair:eq:rounded-level}
 K_u=\{(x,y):e^{\lambda x^2}+e^{\mu y^2}\le e^P+1\}.
\end{equation}
Its intercepts are exactly $\sqrt a$ and $\sqrt b$. Write
\[
 G(s,x,y)=E+F-e^P-1,
 \quad V=EX+FY,\quad
 D_0=e^P P-EX-FY.
\]
All identities below are evaluated on $G=0$ unless indicated otherwise.

\begin{lemma}[Rounded-rectangle phase]
\label{pair:lem:rounded-phase}
For sufficiently small $\eta$ and sufficiently large $K$, the sets
$K_u$ are strictly nested. Their inverse level function $H(x,y)=u$
is smooth outside any fixed initial oval. With
$g=\nabla H$, there is an explicit positive matrix $\mathcal B$ such
that
\begin{equation}\label{pair:eq:rounded-proxy}
 D^2H+g g^{\mathsf T}\preceq\mathcal B,
 \qquad
 \sqrt{\det\mathcal B}\le C\,(xH_x+yH_y)
\end{equation}
at all sufficiently high levels. The constants are uniform in the mesh
parameter. After a fixed smooth flattening at low levels, isotropic
spatial dilation gives a smooth phase $\mathcal F(t,z)$ and a positive
matrix proxy $\mathcal A(t,z)$ with
\begin{equation}\label{pair:eq:rounded-supersolution}
 D^2\mathcal F+\nabla\mathcal F\nabla\mathcal F^{\mathsf T}
 \preceq\mathcal A,
 \qquad -\mathcal F_t\ge\sqrt{\det\mathcal A}.
\end{equation}
Its terminal cost controls both desired complementary Young modulars:
\begin{equation}\label{pair:eq:rounded-terminal}
 \psi_1(|x|/C)+\psi_2(|y|/C)\le C e^{\mathcal F(T,x,y)}.
\end{equation}
The phase diverges at the boundary of its eventual domain.
\end{lemma}

\begin{proof}
First note the entropy identity
\begin{equation}\label{pair:eq:rounded-Q}
 Q:=-G_s=\alpha EX+\beta FY+\gamma D_0.
\end{equation}
Indeed $D\log\lambda=\beta-1=\gamma-\alpha$ and
$D\log\mu=\alpha-1=\gamma-\beta$.
Since $E+F=e^P+1$ and $E,F\ge1$,
\begin{equation}\label{pair:eq:entropy-deficit}
 0\le D_0\le P\min(EX,FY).
\end{equation}
For example, fixing $P$ and taking $E\le F$, the deficit vanishes at
$E=1$ and its derivative with respect to $E$ is $\log(F/E)\le P$.
Thus $D_0\le P(E-1)\le P EX$; exchange the coordinates for the other
bound. Because $|\gamma|P=\Lambda|DT|$ can be made arbitrarily small,
\begin{equation}\label{pair:eq:Q-comparability}
 Q\asymp\alpha EX+\beta FY,
 \qquad Q\ge c\{\min(EX,FY)+\delta V\}>0.
\end{equation}
This proves strict nesting and the existence of the inverse level
function.

Here is a matrix upper bound that does not assume convexity of $H$.
Put $p=\nabla_zG$ and
\[
 \mathcal D=\operatorname{diag}(G_{xx},G_{yy})
 =\operatorname{diag}\{2\lambda E(1+2X),
                         2\mu F(1+2Y)\}.
\]
Implicit differentiation gives
\begin{equation}\label{pair:eq:implicit-Hessian}
 g=\frac uQp,\qquad
 D^2H+gg^{\mathsf T}
 =\frac uQ\mathcal D
   +\frac u{Q^2}(p\otimes G_{sz}+G_{sz}\otimes p)
   +\left\{\frac{u(G_{ss}+Q)}{Q^3}
                                +\frac{u^2}{Q^2}\right\}pp^{\mathsf T}.
\end{equation}
The derivatives that enter this expression obey
\begin{equation}\label{pair:eq:rounded-derivatives}
 G_{sx}=-(\alpha-\gamma)(1+X)p_1,
 \quad G_{sy}=-(\beta-\gamma)(1+Y)p_2,
\end{equation}
\[
 G_{sz}^{\mathsf T}\mathcal D^{-1}G_{sz}\le C P V,
 \qquad |G_{ss}+Q|\le C P V.
\]
For the last bound, differentiate the defining exponentials twice;
$|DX|\le CX$, $|DY|\le CY$, $|DP|+|D^2P|\le C$, and
$V\asymp P e^P$. These estimates also follow directly from
\eqref{pair:eq:rounded-derivatives} and the stated logarithmic derivative
bounds. Cauchy--Schwarz in the metric $\mathcal D$ now proves the
Loewner bound in \eqref{pair:eq:rounded-proxy} with
\begin{equation}\label{pair:eq:explicit-proxy}
 \mathcal B=
 C\frac uQ\mathcal D+
 \left(C\frac{uPV}{Q^3}+\frac{u^2}{Q^2}\right)pp^{\mathsf T}.
\end{equation}
Equivalently, it has the useful diagonal-plus-gradient form
\[
 \mathcal B=\operatorname{diag}(A_1,A_2)+(1+r)gg^{\mathsf T},
 \quad A_i=C\frac uQG_{ii},\quad r=C\frac{PV}{uQ}.
\]
The same estimates with absolute values of the cross and rank-one
coefficients also give $-A-rgg^{\mathsf T}\preceq D^2H\preceq
A+rgg^{\mathsf T}$. Thus the later spatial argument has a bound for
both signs of the Hessian, not just for the exponential Hessian.

We give the determinant bound explicitly. The elementary identities are
\[
 R_0:=p^{\mathsf T}\mathcal D^{-1}p
 =\frac{2XE}{1+2X}+\frac{2YF}{1+2Y}\asymp V/P,
 \qquad xH_x+yH_y=\frac{2uV}{Q}.
\]
The rank-one determinant formula therefore gives
\begin{equation}\label{pair:eq:proxy-ratio}
 \frac{\det\mathcal B}{(xH_x+yH_y)^2}
 \le C\frac{G_{xx}G_{yy}}{V^2}
       \left[1+\left(\frac VQ\right)^2+\frac{uV}{PQ}\right].
\end{equation}
The crucial mixed-exponential estimate is
\begin{equation}\label{pair:eq:mixed-exponential}
 EF(1+X)(1+Y)\le C VQ.
\end{equation}
To check it, take $X\le Y$. Then $Y\asymp P$, $F\asymp e^P$, and
$V\asymp PF$. If $X\ge1$, the left side is at most
$C(EX)(PF)$ and the $\min(EX,FY)$ term in
\eqref{pair:eq:Q-comparability} suffices. If $X\le1$, the left side is
at most $CPF\asymp V$, whereas
$Q\ge c\delta V\ge c\varepsilon_*Pe^{(\Lambda-c_*)T}$.
The last expression is bounded below by a sufficiently large positive
constant after increasing $K$. This proves
\eqref{pair:eq:mixed-exponential}.

Finally
\[
 G_{xx}G_{yy}
 =\frac{4P^2}{ab}EF(1+2X)(1+2Y),\qquad ab=uT,
 \qquad V/Q\le C\delta^{-1}.
\]
Substitution into \eqref{pair:eq:proxy-ratio} gives
\begin{equation}\label{pair:eq:proxy-ratio-final}
 \frac{\det\mathcal B}{(xH_x+yH_y)^2}
 \le C\left[1+\frac{T e^{c_*T}}u\right]\le C.
\end{equation}
The last inequality uses $T=O(\log(e+\log u))$. This proves the asserted
high-level matrix bound.

For clarity, low-level extension costs only fixed constants. Fix a large
$u_0$ and set $H=u_0$ inside $K_{u_0}$. Choose a smooth nondecreasing
convex $\chi$ that vanishes on $(-\infty,u_0]$ and has $\chi'=1$
above $u_0+1$. The function $\chi(H)$ is smooth, including across the
initial oval. A positive proxy for its exponential Hessian is
\[
 \mathcal B_\chi
 =\chi'\,\mathcal B+(\chi''+\chi'^2)gg^{\mathsf T}.
\]
On the fixed transition annulus, all entries and their derivatives are
bounded uniformly in $m$: the allocated widths there have fixed upper
and positive lower bounds, and so do the slope and nesting constants.
For the spatial selector use a fixed positive diagonal majorant $cI$
throughout $H\le u_0+2$, with $c$ large enough to dominate the absolute
Hessian of $\chi(H)$ there. Its gradient rank term has coefficient one.
Blend smoothly to the displayed high-level proxy only on
$u_0+2\le H\le u_0+3$, where $\chi'=1$; the constant $c$ also dominates
on this fixed annulus. Convex combinations preserve the matrix upper
bound. This gives a representation $A+(1+r)\nabla\chi(H)\otimes
\nabla\chi(H)$ with a uniformly positive compact diagonal $A$ and a
bounded $r$. In particular one does not divide $\chi''$ by $\chi'^2$
near the flat region. Above the annulus its determinant has the same
bound as \eqref{pair:eq:proxy-ratio-final}, up to a fixed factor.

Let $s(t)>0$ satisfy $s(t)s'(t)=C_1$, with $s(0)=1$, and put
\[
 \mathcal F(t,z)=\chi(H(z/s(t)))+C_0(T-t).
\]
Choose $C_1$ sufficiently large to dominate the high-level constant in
\eqref{pair:eq:rounded-proxy}; choose $C_0$ to dominate the bounded
transition annulus. Spatial dilation multiplies the positive matrix
proxy by $s(t)^{-2}$, while
\[
 -\mathcal F_t
 =\frac{s'}s\chi'(H)\,z/s\cdot\nabla H+C_0.
\]
This yields \eqref{pair:eq:rounded-supersolution}. A small isotropic
matrix regularization, with a fixed additional time margin, makes the
selector strictly positive even in the flat region.

Every level obeys $|x|\le\sqrt{a(u)}$ and $|y|\le\sqrt{b(u)}$.
The allocation bounds consequently give
$\psi_i(|z_i|/C)\le e^u$ on that level. Fixed spatial dilation and the
identity $\chi(u)=u-O(1)$ prove
\eqref{pair:eq:rounded-terminal}. The finite initial and transition
levels are absorbed in the same fixed constant. As usual, convexity of
$\psi_i$ turns a fixed bound on its modular into a Luxemburg bound after
one further fixed dilation. Strict nesting of compact levels also shows
that the phase diverges at the boundary of their union.
\end{proof}

The formulas also display the quantitative bounds needed for spatial
realization. Differentiating a fixed number of times in logarithmic
level and in the normalized coordinates $\sqrt a\,\partial_x$,
$\sqrt b\,\partial_y$ introduces only fixed powers of $u$ and
$e^{CT}$. The logarithmic growth cap gives polylogarithmic losses in the
level, and the mesh-end cap gives the sharper estimate
\[
 e^{CT_m(\log u)}\le C(m/u)^{C\eta}\qquad(u\le C_1m).
\]
Thus choosing $\eta$ sufficiently small preserves the small conditional
step and oscillation errors in the Gaussian realization.

% END rounded_rectangle_phase.tex


\section{Spatial realization of the phase}

% BEGIN generalized_pair_flow.tex
\subsection{Spatial realization of the two-cost phase}
\label{pair:sec:flow-realization}

We give the spatial argument, including the rank-one Hessian term and
the exact repairs for both Young norms. Write $u=H(z/s(t))$ for the
underlying implicit level and $H_m(t,z)=\chi(u)$ for its flattened
phase. We keep the bounded time-only correction separate, in the
modular below. Above the cutoff annulus $H_m$ differs from $u$ by
one fixed additive constant, so their spatial and time derivatives
agree. All high-level formulas below use this underlying $u$.
Put $v=\nabla H_m$ and $A=\operatorname{diag}(d_1,d_2)$.
The time-dilation margin $\Lambda_{\rm time}$ below is independent
of the already fixed shape coefficient in $P=\Lambda T$.
Set
\[
 B=A+vv^{\mathsf T},\quad K=\sqrt{\det B},\quad
 C=KB^{-1},\quad b_i^2=K/d_i.
\]
These are deliberately the slow curvature widths, rather than the
diagonal entries of the selector. On the whole compact flattened
region choose a fixed sufficiently large diagonal matrix $A=cI$
dominating both signs of its Hessian, and put $r=0$. Blend to the
high-level formula only in a fixed annulus above the point where the
cutoff derivative has become one. All coefficients on this annulus
are uniformly bounded, so this preserves a bounded $r$ and avoids
division by a vanishing cutoff derivative.
The preceding phase calculations supply the following estimates above
the fixed cutoff annulus, with $u\le C_1m$:
\begin{gather}
 -A-rvv^{\mathsf T}\preceq D^2H_m\preceq A+rvv^{\mathsf T},
 \qquad 0\le r\le C,\quad r\le u^{-1}e^{CT_m(\log u)}, \label{pair:eq:engine-hessian}\\
 K\le ue^{CT_m(\log u)},\qquad v_i^2/d_i\le ue^{CT_m(\log u)},
 \qquad -\partial_tH_m\ge c\Lambda_{\rm time} K, \label{pair:eq:engine-phase}\\
 |b_i\partial_i\log f|\le e^{CT_m(\log u)},\qquad
 |\partial_t\log f|\le\Lambda_{\rm time} e^{CT_m(\log u)}
 \quad(f=1+u,d_1,d_2,K,b_1,b_2), \label{pair:eq:engine-scalar-derivatives}\\
 |\partial_t v_i|\le\Lambda_{\rm time} e^{CT_m(\log u)}|v_i|.
 \label{pair:eq:engine-time-gradient}
\end{gather}
The finitely many differentiated comparison estimates used below have
the same form after increasing $C$. All widths, coefficients, and
reciprocals have polynomial bounds in $m$; the flat compact region has
uniform constant versions. The negative time derivative estimate is
only asserted above the cutoff annulus; the separate scalar clock
will cover its complement. For comparisons across the compact region
use the globally smooth quantity $1+H_m$, which is comparable to $u$
at high levels, in place of $1+u$.
The mesh cap ensures, for every fixed
power required in these estimates,
\begin{equation}\label{pair:eq:engine-mesh-floor}
 e^{CT_m(\log u)}\le C'(m/u)^{C'\eta_{\rm path}}
 \qquad(1\le u\le C_1m).
\end{equation}
The path parameter is chosen small after the finite powers are fixed.

Here is a precise order for these parameter choices. First fix the
prefix-data constants, the width regularization exponent, the positive
smoothing kernel and its finitely many derivative orders, the clock
exponent $c_*$, and the shape parameter in $P=\Lambda T$. The
exponential powers denoted $C$ are fixed at this stage: all estimates
are fixed-degree polynomials in $P$, the inverse slope lower bound
$\delta_{\rm slope}^{-1}=\varepsilon_*^{-1}e^{c_*T}$, and the relative
slope derivatives. Dependence of the Hall padding or initial axes on
later parameters enters only multiplicative constants. Choose
$\eta_{\rm path}$ using these fixed exponential powers and
$\eta_{\rm clip}$. Then choose $K$ sufficiently large, of order at
least $\log(1/\eta_{\rm path})$, to suppress the clock convolution
losses, and choose the fixed cutoff $u_0$ large enough to absorb
all remaining compact prefactors. The compact diagonal proxy and scalar
clock may depend on these choices. The metric-memory coefficient,
time-drift margin and finally $\varepsilon$ are chosen afterwards.
This separates exponent choices from compact prefactors and removes
any apparent circularity in the mesh-cap argument.

The slow curvature widths satisfy the exact identities
\begin{equation}\label{pair:eq:canonical-width-identities}
 K^2=d_1d_2+d_1v_2^2+d_2v_1^2,\qquad
 C_{ii}=(d_j+v_j^2)/K\le b_i^2,\qquad d_ib_i^2=K\quad(j\ne i).
\end{equation}
Write $U_*=v_1^2/d_1+v_2^2/d_2$. Then $b_1^2b_2^2=1+U_*$.
For determinant-one $G$ with $G_{ii}\le M_0b_i^2$, put
$\mathcal L(G)=\sqrt{\operatorname{tr}(GC^{-1})}$. Then
\begin{equation}\label{pair:eq:general-lag-size}
 \mathcal L^2\lesssim_{M_0}1+U_*
 \lesssim_{M_0}(1+u)e^{CT_m(\log(1+u))},
 \qquad v^{\mathsf T}Gv\le K\mathcal L^2.
\end{equation}
Indeed $\operatorname{tr}(GA)\le2M_0K$, while
$v^{\mathsf T}Gv\le2M_0K U_*$.

Here is an explicit verification of the scalar derivative bounds,
which is useful near the coordinate axes. In the notation of the
rounded phase let $\mathcal D=\operatorname{diag}(G_{xx},G_{yy})$,
$p=\nabla_zG$, and $R_0=p^{\mathsf T}\mathcal D^{-1}p\asymp V/P$.
Since $d_i=c(u/Q)\mathcal D_{ii}$ for a fixed positive $c$,
\[
 K^2=(u/Q)^2\det\mathcal D\,[c^2+c(u/Q)R_0].
\]
Here $V\asymp Pe^P$ and $R_0\asymp e^P$ on every level. Thus neither
logarithmic differentiation of $R_0$ nor that of the bracket introduces
a small coordinate-gradient denominator. The formulas also give
\[
 e^{-CT_m}\le b_1^2/a,\ b_2^2/b\le e^{CT_m}.
\]
For this estimate use
$\mathcal D_{11}=2(P/a)E(1+2X)$,
$\mathcal D_{22}=2(P/b)F(1+2Y)$, $ab=uT$,
$R_0/Q\ge c/P$, and $R_0/Q\le e^{CT_m}/P$.
In the semiaxis directions $\sqrt a\,\partial_x$ and
$\sqrt b\,\partial_y$, implicit differentiation gives
$|\partial\log u|\le e^{CT_m}$, since $Q\ge c\delta V$.
The identities $Q_s=-G_{ss}$ and $Q_{z_i}=-G_{sz_i}$, together
with $|G_{ss}|\lesssim PV$ and the displayed mixed derivatives of
$G$, bound derivatives of $\log Q$ by $e^{CT_m}$.
The explicit positive formulas for $\mathcal D$, $R_0$, and $K^2$
then prove \eqref{pair:eq:engine-scalar-derivatives}.
Convert to the slow widths by their comparison with the semiaxes.
Repeated differentiation introduces only fixed powers of
$P,\delta_{\rm slope}^{-1}$ and their controlled logarithmic derivatives, so
the same assertion holds through every fixed required order.
For the time derivatives use isotropic dilation and
$u_t/u=O(\Lambda_{\rm time} V/Q)$; each nonzero $v_i$ is a coordinate $z_i$
times a positive coefficient with logarithmic derivative bounded
by $\Lambda_{\rm time} e^{CT_m}$. A zero coordinate remains zero under
dilation. This proves \eqref{pair:eq:engine-time-gradient}.

Use the rectangular factor $S_C$ with rows
$\sqrt{d_1/K}\,e_1^{\mathsf T}$,
$\sqrt{d_2/K}\,e_2^{\mathsf T}$, and $v^{\mathsf T}/\sqrt K$.
Thus $S_C^{\mathsf T}S_C=C^{-1}$ and
$\mathcal L=\|S_CG^{1/2}\|_{\rm HS}$. The first two rows after
multiplication by $G^{1/2}$ are bounded, by
\eqref{pair:eq:canonical-width-identities}. Their derivatives use
\eqref{pair:eq:engine-scalar-derivatives}. For the last row use
\eqref{pair:eq:engine-hessian}: with $D=A+rvv^{\mathsf T}$,
$|D^2H_m(w,z)|\le\sqrt{D(w,w)D(z,z)}$. Consequently
\begin{align}\label{pair:eq:general-lag-derivatives}
 \|b_i\partial_i(S_CG^{1/2})\|_{\rm HS}
 &\lesssim_{M_0}\sqrt{1+u}\,e^{CT_m(\log(1+u))},\notag\\
 \|\partial_t(S_CG^{1/2})\|_{\rm HS}
 &\lesssim_{M_0}\Lambda_{\rm time}\sqrt{1+u}\,e^{CT_m(\log(1+u))}.
\end{align}
For example the square of the differentiated-gradient contribution is
at most $K^{-1}D(b_ie_i,b_ie_i)\operatorname{tr}(GD)$, and both
factors are at most $K$ times a fixed power of $e^{CT_m}$.
The time estimate uses \eqref{pair:eq:engine-time-gradient}.
This factor calculation avoids small eigenvalue denominators.

\begin{proposition}[Two-cost spatial realization]\label{pair:prop:flow}
The phase with these quantitative estimates yields
\[
 \E\inf_{\operatorname{div}J=P_hW^0,\ J\cdot n=0}
 \max_{i=1,2}\|J_i\|_{L^{\psi_i}}
 \le C\sqrt{\log(1/h)}.
\]
\end{proposition}
\begin{proof}
Put $m=\log(1/h)$ and fix $0<\eta_{\rm clip}<1/10$.
Choose $\eta_{\rm path}$ so small that all exponents
$C\eta_{\rm path}$ below are less than $\eta_{\rm clip}$ and $1/4$.
Set
\[
 \delta=\varepsilon^2/m,\qquad E_0=m^{1/2-\eta_{\rm clip}},
 \qquad s=hm^{C_s},
\]
with $C_s$ large fixed. Write $t_k=k\delta$ and start from $U_0=0,G_0=I$.
At an active parent $Q$ choose an old-measurable interior representative
$x_Q$ and set $C_Q=C(t_k,U_k(x_Q))$. If
\[
 G^{-1/2}C_QG^{-1/2}
 =O\operatorname{diag}(\lambda,\lambda^{-1})O^{\mathsf T},
 \qquad
 \lambda^\#=\min\{K_0,\max\{K_0^{-1},\lambda\}\},
\]
put
\[
 G_+=G^{1/2}O\operatorname{diag}
       (\lambda^\#,(\lambda^\#)^{-1})O^{\mathsf T}G^{1/2}.
\]
Thus the two reciprocal relative eigenvalues are clipped reciprocally.
The scalar inequalities
$K_0^{-1}\le\lambda^\#\le K_0$ and
$\lambda^\#\le K_0^{-1}+\lambda$ give
\[
 \det G_+=1,\quad K_0^{-1}G\preceq G_+\preceq K_0G,\qquad
 G_+\preceq K_0^{-1}G+C_Q.
\]
Take the ordinary Whitney split in this chart down to minimum leg $s$.
Apply its local ray inverse, multiply the current by $\sqrt\delta$,
and use even coarse clipping at $E_0$. Write $\Delta U_j$ for the
retained increment born at generation $j$, in metric
$G_j=A_jA_j^{\mathsf T}$, and $\ell_j$ for its child leg.
On that birth child its envelope $E_j$ satisfies
\[
 \|A_j^{-1}\Delta U_j\|_\infty\le C\sqrt\delta E_j,\qquad
 |A_j^{-1}(\Delta U_j(x)-\Delta U_j(x'))|
 \le C\sqrt\delta E_j
       \left(\frac{|A_j^{-1}(x-x')|}{\ell_j}\right)^\zeta ,
\]
for one fixed $0<\zeta<1$. Each retained $E_j$ is at most $CE_0$
and has a uniform conditional Gaussian tail.
Freeze geometric terminal cells. The depth is at most $C_*m$;
take horizon $T=C_*m\delta$ and eventually choose $\varepsilon$ so
all phase time restrictions hold. Stop the whole modular at its first
crossing of a large fixed constant.

Here is the full geometric bootstrap. Inductively maintain the
pointwise diagonal cap $G_{k,ii}\le M_0b_i(t_k,U_k)^2$, validity of
the phase, and per-generation width ratios in $[1-c_0,1+c_0]$,
for a fixed sufficiently small $c_0>0$. These hold initially.
If a current parent is a descendant of a birth child at generation
$j\le k$, its diameter divided by that birth leg, in the birth
metric, is at most $C\vartheta^{k-j}$.
Indeed consecutive charts have relative distortion at most
$\sqrt{K_0}$; choose the intrinsic Whitney contraction so small
that its product with this fixed distortion is $\vartheta<1$.

For $x,x_0$ in the same current parent, use the birth cap at $x_0$
and the width-ratio induction along that fixed spatial point:
\[
 \begin{split}
 |\Delta U_{j,i}(x)-\Delta U_{j,i}(x_0)|
 &\le C\sqrt\delta\,\sqrt{G_{j,ii}}\,E_j(x_0)
                     \vartheta^{\zeta(k-j)}\\
 &\le C\sqrt\delta\,b_i(t_k,U_k(x_0))
       \bigl[(1-c_0)^{-1}\vartheta^\zeta\bigr]^{k-j}E_j(x_0).
 \end{split}
\]
Set $r_{\rm osc}=(1-c_0)^{-1}\vartheta^\zeta<1$ and fix
$r_{\rm osc}^2<\rho<1$. Define, pointwise along the actual birth
history,
\[
 \mathcal T_k(x)=\sum_{j\le k}\rho^{k-j}E_j(x)^2,
 \qquad \mathcal T_{k+1}=\rho\mathcal T_k+E_{k+1}^2.
\]
Thus $\mathcal T_k\le CE_0^2$. Summing the displayed birth estimate
and using
$\sum_{\ell\ge0}(r_{\rm osc}^2/\rho)^\ell<\infty$ gives
\[
 |U_{k,i}(x)-U_{k,i}(x_0)|
 \le C\sqrt\delta\,b_i(t_k,U_k(x_0))
                         \sqrt{\mathcal T_k(x_0)}.                 \tag{*}
\]
This deterministic estimate is anchored at $x_0$; it does not
presuppose any bound for the phase at the other point.

Before a modular stop each active cell has a witness $x_0$ with
$1+H_m(t_k,U_k(x_0))\le C_1m$, because its area is at least $cs^2$
and its integral of $e^{H_m}$ is bounded by the modular cap.
The witness need not be chosen predictably: it is used only in the
deterministic bootstrap. Integrate the coefficient derivative bounds
along the straight state segment from $U_k(x_0)$ to $U_k(x)$.
Up to the first possible failure of factor-two bounds for
$1+H_m,d_i,K,b_i$, estimate (*) gives relative logarithmic change
at most
\[
 C\sqrt\delta E_0e^{CT_m}
 \le C\varepsilon m^{-\eta_{\rm clip}+C\eta_{\rm path}}=o(1).
\]
It prevents that first failure and hence continues to every $x$.
A state segment cannot escape the phase domain while this comparison
holds: properness of the nested levels would force $H_m$ to tend to
infinity. Thus all old states have $H_m\le 2C_1m$ and their widths
and coefficients are relatively $1+o(1)$ throughout the parent.

At the representative the new metric satisfies
$G_{+,ii}\le(K_0^{-1}M_0+1)b_i^2$. Choose $M_0$ so large that
$K_0^{-1}M_0+1\le(1-2c_1)M_0$ for some $c_1>0$.
The preceding representative comparison gives the same bound with
margin $c_1$ at every old point. The fresh displacement divided by
its old widths is at most $C\sqrt{M_0\delta}E_0$.
The same first-failure argument along this fresh state segment and
the short time interval $\delta$ gives relative coefficient change
$o(1)$ and an upper level bound $C_1'm$. The new diagonal cap follows
from its strict old margin. In particular the new width ratios lie
in $[1-c_0,1+c_0]$ for large $m$, closing the bootstrap used in (*).
The coefficient bounds now give polynomial chart distortion.

For completeness the lag contraction is also explicit. If $e^d$
is the larger relative eigenvalue of $G$ against a fixed $C_Q$,
then $\mathcal L(G)^2=2\cosh d$. Clipping replaces $d$ by
$(d-\log K_0)_+$, so
$\mathcal L(G_+)^2\le K_0^{-1}\mathcal L(G)^2+2$.
Representative and fresh-state comparisons cost, by
\eqref{pair:eq:general-lag-derivatives}, at most
\[
 C\sqrt{\delta(1+u)}e^{CT_m}(\sqrt{\mathcal T}+E)
 +C\Lambda_{\rm time}\delta\sqrt{1+u}e^{CT_m}
 \le C\varepsilon(\sqrt{\mathcal T}+E)+C\Lambda_{\rm time}\varepsilon^2.
\]
Consequently, with $q_1<1$ fixed,
\[
 \mathcal L'\le q_1\mathcal L+
 C+C\varepsilon(\sqrt{\mathcal T}+E)
                  +C\Lambda_{\rm time}\varepsilon^2.
\]
Choose $\varepsilon$ after $\Lambda_{\rm time}$. Squaring with a
fixed slack in $q_1<1$ gives
$\mathcal L'^2\le q_2\mathcal L^2+C+
C\varepsilon^2(\mathcal T+E^2)$ with $q_2<1$.
Adding $\mathcal T'=\rho\mathcal T+E^2$ and taking $\varepsilon$
small gives,
for $S=\mathcal L^2+\mathcal T$ and $q<1$,
\begin{equation}\label{pair:eq:general-memory}
 S'\le qS+C(1+E^2),\qquad K'S'-KS\le-cKS+CK(1+E^2).
\end{equation}
The relative change of $K$ is uniformly $o(1)$.

Use $\Xi=H_m(t,U)+C_0(T-t)+\kappa\delta KS$.
The predictable linear tilt has Gaussian variance at most
$C\delta K\mathcal L_+^2$; even scalar clipping decreases its
exponential moment. The rank-one Hessian term requires a scalar
comparison; the gradient itself is not assumed comparable along a
segment. Write $\xi$ for the retained increment and
$f(s)=H_m(t,U+s\xi)$. The proved comparisons of $A,K,b_i$ give
\[
 f''(s)\le a+\widehat r f'(s)^2,\qquad
 a=C\delta KE^2,\qquad
 \widehat r=C(1+u)^{-1}e^{CT_m(\log(1+u))}.
\]
Here $\xi^{\mathsf T}A\xi\le a$ follows from the diagonal cap and
$d_ib_i^2=K$. Moreover
$|f'(0)|\le C\sqrt{\delta K}\,\mathcal L_+E$.
The mesh cap and $E\le CE_0$ imply
$\widehat r|f'(0)|+\widehat r a=o(1)$:
the square of the first term and the second term are bounded by
a fixed power of $e^{CT_m}$ times $C\delta E_0^2=o(1)$.
Scalar Riccati comparison therefore gives
\[
 f(1)-f(0)\le f'(0)+Ca+C\widehat r f'(0)^2.
\]
Indeed compare $z=\widehat r f'$ to
$z'=\widehat r a+z^2$; small initial value and forcing bound its
increment by $C(\widehat r a+z(0)^2)$. The case $\widehat r=0$
is immediate. The spatial Taylor contribution is thus at most
$C\delta K(1+\widehat r\mathcal L_+^2)E^2$.
Use $\widehat r$ in place of $r$ in the following estimates;
it has the same uniform and mesh-floor bounds.

The moment bounds remain uniform. Besides $\delta K\lesssim\varepsilon^2$,
the mesh cap yields
\[
 \delta rK\mathcal L^2\lesssim\delta u e^{CT_m}\lesssim\varepsilon^2,
 \qquad
 \delta rK\mathcal T\lesssim
 \varepsilon^2m^{-2\eta_{\rm clip}+C\eta_{\rm path}}
 \lesssim\varepsilon^2.
\]
Thus the conditional envelope moment and Cauchy--Schwarz bound the
positive quadratic contribution by $C\delta K+C\delta rKS$
in the logarithmic moment. The phase gives negative time drift
$-c\Lambda_{\rm time}\delta K$; the scalar clock covers its compact flat
and transition regions. Together with \eqref{pair:eq:general-memory},
the logarithmic moment of $\Xi'-\Xi$ is at most
\[
 -c\Lambda_{\rm time}\delta K-c\kappa\delta KS
 +C\delta K\mathcal L^2+C\delta rKS+C_\kappa\delta K.
\]
Since $r$ is globally bounded, choose $\kappa$ to pay the two
terms proportional to $S$, then $\Lambda_{\rm time},C_0$ to pay the constant,
and $\varepsilon$ last. Spatial integration with prescribed freezing
gives a nonnegative supermartingale with bounded constant initial
value. Its maximal inequality gives a fixed positive probability
of completing without a modular stop. The final modular and the
phase's level-set inclusion imply
$\|U_i\|_{L^{\psi_i}}\le C$ for both $i$.

We restore the exact divergence using both suprema in the sparse
repair proof. With $\sigma=\sqrt\delta$,
$T_0=\sigma\sqrt m=\varepsilon$, and $D_0=m/E_0^2=m^{2\eta_{\rm clip}}$,
its support has area at most $e^{-cE_0^2}$ and
\[
 \|J_1\|_\infty\le CT_0D_0B,\qquad
 \|J_2\|_\infty\le CT_0D_0/B.
\]
Allow $B<1$ by exchanging the axes. Put
$R_i=\rho_i(e^{c'E_0^2})$ for small fixed $c'>0$ and choose
\[
 B=\sqrt{\widehat R_1/\widehat R_2},\qquad
 \widehat R_i=\min\{R_i,m\}.
\]
The prefix condition implies $R_1R_2\gtrsim E_0$.
To recall the short proof, if
$p_j=\rho_1(e^{2^j})^2\rho_2(e^{2^j})^2\,2^{-j}$
tended to zero, then for all $j-r\le k\le j$ monotonicity would give
$\rho_1(e^{2^k})^2\rho_2(e^{2^j})^2\,2^{-k}\le2^rp_j$.
Taking $r$ proportional to $\log(1/p_j)$ produces arbitrarily many
terms bounded below in the prefix sum, a contradiction. There are
enough preceding indices because both inverse widths have fixed
positive lower bounds. Thus $\inf_jp_j>0$; interpolation between
dyadic levels gives $\rho_1(e^v)\rho_2(e^v)\gtrsim\sqrt v$.
Each $R_i$ is bounded below by a fixed positive constant, so truncating
at $m$ still leaves
$\widehat R_1\widehat R_2\gtrsim E_0$, since $E_0\le\sqrt m$.
The ratio $B$ is polynomially bounded in both directions.
A field of supremum $L$
on area $p$ has norm at most $L/\rho_i(1/p)$. Both normalized repair
norms are therefore at most
\[
 C T_0D_0/\sqrt{E_0}
 =O(m^{(5/2)\eta_{\rm clip}-1/4})=o(1).
\]
This argument also covers costs stronger than Gaussian; it does
not use a false $\psi_2$ embedding.

Here are the exact terminal corrections, with normalization made
explicit. The preceding retained currents and clipping repairs have
divergence $\sqrt\delta f_{\rm leaf}$. Thus $f_{\rm leaf}$ denotes
the \emph{raw} revealed Gaussian density, before multiplication by
$\sqrt\delta$. At each generation there are at most $Cs^{-2}$
attempted cells, because every child has area at least $cs^2$.
The number over all generations is at most $Cm s^{-2}=\exp(O(m))$.
Every fresh cell density has conditional standard deviation at
most $C/s$. Gaussian tails with threshold $C_A\sqrt m/s$, followed
by the union bound, give simultaneous density bounds with failure
at most $e^{-Am}$, where $A$ is arbitrarily large by increasing
$C_A$. Summing at most $C_*m$ ancestral contrasts proves
\[
 \|f_{\rm leaf}\|_\infty\le Cm^{3/2}/s
\]
with probability $1-o(1)$. This estimate concerns the full revealed
contrasts: the sparse repairs have already restored the divergences
removed by clipping.

The raw correction $(P_h-I)f_{\rm leaf}$ has mean zero in every
$h$-square. Row primitives after subtracting row means, followed
by the vertical primitive of those means, give zero-normal
currents in that square of both suprema at most
$Ch\|f_{\rm leaf}\|_\infty=O(m^{3/2-C_s})=o(1)$.
The corresponding normalized correction is smaller by
$\sqrt\delta$, although we only need the displayed raw bound.

By the geometric stopping rule and polynomial chart cap, all
terminal leaves have physical diameter at most $h m^{C_2}$ for
a fixed $C_2$. Choose an $h$-aligned dyadic length $D$ satisfying
\[
 4(\max_Q\operatorname{diam}Q+\sqrt2h)\le D
 \le C h m^{C_2};
\]
for large $m$ it is less than $1/4$. Partition the unit square
into dyadic $D$-squares and assign each terminal leaf to the
square containing one of its representatives. Its final-grid
projection is supported in that square and its immediate
neighbors. Enclose these supports in an $h$-aligned square of
side $4D$; at the domain boundary slide that square inside the
unit square. It can contain every assigned support, because their
coordinate spans are at most $3D$. Each enlargement lies within
$O(D)$ of its original grouping square, so the enlarged boxes
have a universal bounded overlap.

Condition on the full revealed history. For a grouping box $R$,
the unresolved source is the sum of the bridges on its assigned
leaves. These are independent isonormal processes on the
orthogonal spaces $L^2_0(Q)$ of the assigned leaves. If
$\Pi_R$ is their orthogonal-sum projection, its raw projected
cell masses divided by $h$ are
\[
 g_v=W\bigl(\Pi_R(\mathbf1_v/h)\bigr).
\]
The tests $\mathbf1_v/h$ are orthonormal, and $\Pi_R$ is a
contraction. Thus the conditional covariance of $(g_v)$ is at
most the identity. Its sum is exactly zero: the grouping box
contains every assigned leaf, and each leaf bridge annihilates
its constant. The capacity lemma applies to the $4D/h$ square
mesh. Its mean is $O(\log(D/h))=O(\log m)$; concentration with
deviation $C_A\sqrt m$ gives failure $e^{-Am}$ for arbitrary
fixed $A$, after choosing $C_A$. There are at most $h^{-2}=e^{2m}$
groups, so the union bound gives simultaneous raw capacity
$O(\sqrt m)$ with probability $1-o(1)$. Face interpolation gives
divergence $g_v/h$, exactly the required projected bridge density,
and zero normal trace on each enlarged box. Bounded overlap
then gives a sum with both raw suprema $O(\sqrt m)$.
These bounds control both Young norms after a fixed dilation.
Their normalized costs would be $O(\sqrt{\delta m})=O(\varepsilon)$.

Undo $\sqrt\delta$ and add every repair. The divergence is exactly
$P_hW^0$ and the desired gauges hold with fixed positive probability.
For clarity, the final expectation upgrade is also quantitative.
For a mean-zero cell vector $g$, the grid cut formula and
$|\partial S|\ge c\sqrt{\min(|S|,|S^c|)}$ give capacity at most
$C\|g\|_2$ by Cauchy--Schwarz. Face interpolation produces a
current of divergence $h^{-1}g$ and both suprema $C\|g\|_2$.
A constant of size $L$ has $L^{\psi_i}$ norm
$L/\rho_i(1)$, so the minimum-flow gauge $\mathcal F_h$ satisfies
\[
 |\mathcal F_h(g)-\mathcal F_h(g')|
 \le C_{\psi_1,\psi_2}\|g-g'\|_2 .
\]
The centering projection is a Euclidean contraction, and hence
this Lipschitz constant is independent of the mesh for standard
Gaussian histogram coefficients. The preceding construction gives
$\PP(\mathcal F_h\le C\sqrt m)\ge\varepsilon_0>0$ uniformly in $m$.
The lower Gaussian concentration bound therefore implies
\[
 \E\mathcal F_h\le C\sqrt m+
 C_{\psi_1,\psi_2}\sqrt{2\log(1/\varepsilon_0)}.
\]
This is the claimed estimate. The finitely many small values of
$m$ are absorbed by enlarging its constant.
\end{proof}

% END generalized_pair_flow.tex


\section{Necessity and exact examples}

% BEGIN linear_stage_input.tex
\subsection{The entropy obstruction}

We give the lower argument for the coordinate cost
\(\Psi(x,y)=\theta(|x|)+|y|\); exchanging the axes gives the convention
used in the question.  Here \(\theta\) is finite and convex on
\([0,\infty)\), \(\theta(0)=0\), and \(\theta(t)\ge t\).  In particular
its inverse is continuous.  Put
\[
 a_j=\max\left\{1,
 \frac{e^{2^j}}{2^{j/2}\theta^{-1}(e^{2^j})}\right\}.
 \tag{L1}\label{lower:eq:aj}
\]

\begin{lemma}[Entropy criterion]\label{lower:lem:entropy}
For a sequence \(a_j\ge1\),
\[
 \sup_{\substack{q_j\ge0\\\sum_jq_j\le1}}
       \sum_j\frac{q_j\sqrt{\log(1/q_j)}}{a_j}<\infty
 \quad\Longleftrightarrow\quad
 \sum_j e^{-c a_j^2}<\infty\ \hbox{for some }c>0.
\]
The supremum may be restricted to finitely supported sequences.
\end{lemma}
\begin{proof}
For sufficiency choose \(A\) large and split each summand according as
\(\sqrt{\log(1/q_j)}\le Aa_j\) or not.  The first sum is at most
\(A\).  In the second, writing \(x=\sqrt{\log(1/q_j)}>Aa_j\),
\(e^{-x^2}x/a_j\le C_Ae^{-A^2a_j^2/2}\), which is summable for
sufficiently large \(A\).

Conversely suppose every such exponential series diverges.  Fix
\(0<\gamma<1\), discard any finite initial segment, and, for each large
\(A\), select a finite set of remaining indices with
\[
 1\le Z_A:=\sum_j e^{-Aa_j^2}\le2,
 \qquad q_j=\gamma e^{-Aa_j^2}/Z_A.
 \tag{L2}\label{lower:eq:gibbs}
\]
This is possible by taking a partial sum up to its first crossing of
one; each term is at most \(e^{-A}\).  Then \(\sum q_j=\gamma\) and
\(\log(1/q_j)\ge Aa_j^2\), so the entropy sum is at least
\(\gamma\sqrt A\).
\end{proof}

\begin{lemma}[A bounded-cost Brownian control]\label{lower:lem:control}
If the series in Lemma~\ref{lower:lem:entropy}, with
\eqref{lower:eq:aj}, diverges for every \(c>0\), then for every
\(D>0\) there is a bounded predictable upper-triangular matrix process
\(L_t\), with nonnegative diagonal entries, such that its martingale
\(Z_t=\int_0^tL_s\,dB_s\) stays in a fixed compact set and
\[
 \E\Psi(Z_1)\le3,
 \qquad \E\int_0^1\sqrt{\det L_t}\,dt\ge D.
 \tag{L3}\label{lower:eq:control}
\]
The bounds on the process and its compact set may depend on \(D\).
\end{lemma}
\begin{proof}
We first construct an uncapped process and subsequently stop it in
the state variable.  Set
\[
 S_j=2^j,\quad H_j=e^{8S_j},\quad
 \widetilde a_j=\max\{1,H_j/(\theta^{-1}(H_j)\sqrt{S_j})\}.
\]
Since the quantity before taking the maximum is \(\sqrt8\) times
the corresponding quantity at index \(j+3\) in
\eqref{lower:eq:aj}, the exponential series for
\(\widetilde a_j\) also diverges for every constant.  Choose the
finite probabilities \(q_j\) by \eqref{lower:eq:gibbs}, using these
\(\widetilde a_j\) and only indices \(j\ge j_0\), where \(j_0\)
will be fixed below.  Take \(\gamma\) to be a sufficiently small
absolute constant.  Define
\[
 b_j=(\sqrt{S_j}\widetilde a_j)^{-1},\quad
 \ell_j=q_jb_j,\quad K_j=\log(1+q_j^{-2}),\quad
 w_j(v)=\frac{S_j}{K_j(v+\ell_j^2)}.
 \tag{L4}\label{lower:eq:stage}
\]

In operational stage \(j\), start a fresh real Brownian motion
\(U\) at zero and stop at
\[
 T_j=\tau_{\ell_j}\wedge\tau_{\{\pm b_j\}}\wedge b_j^2,
 \qquad R_j=\int_0^{T_j}\sqrt{w_j(v)}\,dv,
\]
where \(\tau_\ell\) is the first time Tanaka-normalized local time
at zero reaches \(\ell\).  Success means that \(\tau_{\ell_j}\)
occurs first; at success \(U_{T_j}=0\).  The following estimates
have absolute constants, uniformly in all the parameters:
\[
 \PP(\text{failure})\le Cq_j,\qquad
 \E[R_j;\text{success}]\ge
       \frac{c q_j\sqrt{K_j}}{\widetilde a_j},\qquad
 0\le R_j\le\frac2{\widetilde a_j\sqrt{K_j}}.
 \tag{L5}\label{lower:eq:resources}
\]
For completeness, Tanaka's formula itself follows by applying It\^o's
formula to smooth convex approximations of $|x|$. It reads
$|U_t|=M_t+L_t^0$, where $M_t=\int_0^t\operatorname{sgn}(U_s)dU_s$
is Brownian motion because its quadratic variation is $t$. The regulator
$L^0$ increases only when $|U|=0$, so the elementary reflection
identity gives $L_t^0=-\min_{s\le t}M_s$. Reflection at the first
crossing of a positive level yields
$\PP(\tau_\ell>v)=2\Phi(\ell/\sqrt v)-1\asymp
\min(1,\ell/\sqrt v)$.
At the bounded stage stopping time, Tanaka's formula gives
\(\E|U_{T_j}|=\E L^0_{T_j}\le\ell_j\), so exit at
\(\pm b_j\) has probability at most \(q_j\).  The reflection
principle and L\'evy's local-time identity give
\(\PP(\tau_\ell>v)\asymp\min(1,\ell/\sqrt v)\); the time-cap
failure probability is therefore \(O(q_j)\).  Subtracting the total
failure probability shows, for a sufficiently small absolute
\(c_0>0\),
\[
 \PP(T_j>v,\text{success})\ge c\ell_j/\sqrt v,
 \qquad \ell_j^2\le v\le c_0b_j^2.
\]
Integration against \(\sqrt{w_j(v)}\,dv\) proves the middle
bound of \eqref{lower:eq:resources}; direct integration proves the
last one.

Conditioned on all stages succeeding, their resources are independent.
Their sum has mean \(\mu_A\ge c\gamma\sqrt A\).  Moreover
\(\log(1/q_j)\ge A\widetilde a_j^2\), so its variance is at most
\(C\mu_A/\sqrt A\).  The probability of all successes is at least
\(1-C\gamma\).  Consequently, for large \(A\), with probability
at least \(1/2\), all stages succeed and their total resource exceeds
\(R:=\mu_A/4\), where \(R\to\infty\).

Run these stages using the first coordinate \(B^1\) of a planar
Brownian motion, with ordinary time given by
\[
 dt=\sqrt{w_j(v)}\,dv/R,\qquad
 \sigma_t^2dt=dv,\quad \beta_t^2dt=w_j(v)dv.
\]
Stop at the first failed stage, the end of the finite stage list, or
ordinary time one.  Set \(\sigma=\beta=0\) afterward.  During activity
\(\sigma\beta=R\).  The stage clocks satisfy
\[
 \int_{\text{stage }j}\beta_t^2dt\le S_j,
 \qquad s_t:=\int_0^t\beta_r^2dr\le2S_j
       \quad\hbox{while visiting stage }j.
 \tag{L6}\label{lower:eq:clock}
\]
The coefficients are bounded for this fixed finite list and depend
only on \(B^1\).  The process \(M=\int\sigma\,dB^1\) is confined
to \([-b_j,b_j]\) during stage \(j\), and restarts from zero after
every success.  Define
\[
 Y_t=\exp\left(\int_0^t\beta_r\,dB_r^2-s_t/2\right),
 \qquad X_t=Y_tM_t.
\]
Then the martingale \((X,Y-1)\) has integrand
\[
 \mathcal L_t=
 \begin{pmatrix}Y_t\sigma_t&M_tY_t\beta_t\\0&Y_t\beta_t\end{pmatrix},
 \qquad \sqrt{\det\mathcal L_t}=Y_t\sqrt R
       \quad\hbox{during activity}.
 \tag{L7}\label{lower:eq:triangular}
\]

Now stop at the first crossing \(Y=H_j\) during stage \(j\), calling
this time \(\tau\).  To estimate its cost in reward, use the
\emph{uncapped} probability measure \(dQ=Y_1\,d\PP\).
Conditionally on the complete \(B^1\) path, the coefficients and
clock are fixed, and under \(Q\), \(\log Y\) is a Brownian motion
plus half its variance clock.  Also the \(B^1\) marginal is unchanged,
since \(\E(Y_1\mid B^1)=1\).  By \eqref{lower:eq:clock} and reflection,
\[
 Q(\tau\le1)\le\sum_{j\ge j_0}2e^{-49S_j/4}.
\]
Choose \(j_0\) so that this sum is at most \(1/8\).  The event of
all successes and total resource at least \(R\) keeps the uncapped
process active throughout \([0,1]\).  Therefore
\[
 \E\int_0^{1\wedge\tau}\sqrt{\det\mathcal L_t}\,dt
 =\sqrt R\int_0^1Q(t<\tau,\text{active at }t)\,dt
 \ge(1/2-1/8)\sqrt R.
\]
For the terminal cost, let \(j\) be the last visited stage.  The
stopped variables satisfy \(|M|\le b_j\), \(Y\le H_j\), and
\(b_j\le\theta^{-1}(H_j)/H_j\).  Convexity gives, pathwise,
\[
 \theta(|X|)\le
 \theta\big((Y/H_j)\theta^{-1}(H_j)\big)\le Y.
\]
The stopped \(Y\) is a bounded positive martingale of mean one, so
\(\E\theta(|X|)\le1\) and \(\E|Y-1|\le2\).  Its integrand and
entire state path are bounded, since only finitely many stages occur.
Taking \(L=\mathcal L\mathbf1_{t<\tau}\) and then \(A\) sufficiently
large proves \eqref{lower:eq:control}.
\end{proof}


% END linear_stage_input.tex


% BEGIN bubble_template.tex
\subsection{A concrete affine-orthogonal bubble}
\begin{lemma}\label{pair:lem:template}
On $T=\{(x,y):x,y\ge0,\ x+y\le1\}$ there is a nonzero continuous
piecewise affine function $b$, compactly supported in its interior, such
that $\int_T b\ell=0$ for every affine function $\ell$. For a uniform
point $U$ of $T$, the law of $c_b^{-1}\nabla b(U)$ is bounded, centered,
has covariance $I$, and is invariant under all signed coordinate
permutations. Moreover $b$ is affine on a fixed subdivision into congruent
right-isosceles triangles of a dyadic leg $2^{-k_0}$.
\end{lemma}
\begin{proof}
Take four centres $c_{ij}=(1/8+i/8,1/8+j/8)$, $i,j\in\{0,1\}$,
and radius $r=1/64$. The four squares of radius $r$ about these centres
are disjoint and lie strictly inside $T$. Put
\[
 b(z)=\sum_{i,j=0}^1(-1)^{i+j}(r-\|z-c_{ij}\|_\infty)_+.
\]
All four tents have equal mass. Their first moments are that mass times
their centres. The checkerboard signs cancel the total mass and both
first moments. Each tent's gradient distribution is supported on
$\{0,\pm e_1,\pm e_2\}$ and is invariant under signed coordinate
permutations; changing the sign of a tent preserves this law. Thus its
covariance is $c_b^2I$ for a fixed $c_b>0$, and its mean is zero.
All break lines have dyadic coordinates and directions horizontal,
vertical, or slope $\pm1$. Choose a sufficiently fine dyadic square
lattice through all centres and tent boundaries. In each small square
choose the diagonal that contains its break line, if any; otherwise
choose either diagonal. The resulting triangles, including the halves
along the diagonal boundary of $T$, are congruent right-isosceles
triangles, and $b$ is affine on each. This proves every assertion.
\end{proof}

% END bubble_template.tex


% BEGIN bounded_bubble_realization.tex
\subsection{Bounded-state spatial realization}

Here $\Psi:\mathbb R^2\to[0,\infty)$ is any finite convex function with $\Psi(0)=0$. We record carefully the bounded-state extension of the Gaussian bubble
construction described below.  It prevents any assumption about Gaussian
integrability of \(\theta\) from entering the argument.  Let
\(D=(0,1)^2\), let \(W\) be white noise on \(D\), and let \(P_m\)
be averaging on its dyadic squares of side \(2^{-m}\).

\begin{lemma}[Bounded-state bubble realization]\label{lower:lem:spatial}
Let \(L\) be a bounded predictable upper-triangular process with
nonnegative diagonal entries, and suppose
\(Z_t=\int_0^tL_s\,dB_s\) stays in a fixed compact set.  For every
\(\delta>0\) there are random mean-zero Lipschitz potentials \(u_m\)
such that, for every realization of their Gaussian signs,
\begin{align}
 \int_D\Psi(\nabla u_m)&\le\E\Psi(Z_1)+\delta+o(1),
       \label{lower:eq:pathwise-cost}\\
 \liminf_{m\to\infty}\frac{\E W(P_mu_m)}{\sqrt m}
   &\ge c_*\left(\E\int_0^1\sqrt{\det L_t}\,dt-\delta\right).
       \label{lower:eq:spatial-reward}
\end{align}
Here \(c_*>0\) is universal.  The error and the uniform bound on
\(\nabla u_m\) may depend on the chosen control and on \(\delta\).
\end{lemma}
\begin{proof}
\emph{Compact-state approximation.}
Suppose \(|Z_t|\le H\).  Approximate \(L\) in
\(L^2(\PP\otimes dt)\) by uniformly bounded upper-triangular simple
processes \(V\), constant on a finite deterministic partition and
with block matrices bounded continuous functions of finitely many
earlier Brownian increments.  Density of predictable simple processes,
followed by finite-dimensional continuous approximation, gives this
assertion.  Clip their diagonal entries to a fixed nonnegative interval,
clip the off-diagonal entry, and add \(\varepsilon I\).  The approximation
still converges, and each fixed \(V\) has compact, uniformly invertible
matrix range.

Write \(Z^V=\int V\,dB\), and stop it at its first exit \(\tau_V\)
from the ball of radius \(H+1\).  Doob's inequality gives
\[
 \PP(\tau_V\le1)\le4\E\int_0^1|V_t-L_t|^2dt.
\]
The stopped states stay in this fixed ball.  Uniform continuity of
\(\Psi\) on the ball and convergence in probability give convergence
of their terminal expected costs to \(\E\Psi(Z_1)\).  The rewards
also converge: \(\sqrt{\det V}\) converges in mean by uniform
continuity on the bounded matrix range, and stopping loses at most a
constant times \(\PP(\tau_V\le1)\).  Fix a sufficiently accurate
\(V\) so both errors are less than \(\delta\).  No cost outside a
compact state set has been used.

\emph{Bounded innovations and stopping.}
Use the reference piecewise affine bubble \(b\) on a right-isosceles
triangle \(T\) from Lemma~\ref{pair:lem:template}.  It is compactly supported in
\(T\), satisfies \(\int_Tb\ell=0\) for every affine \(\ell\), and,
for a uniform point \(U\in T\),
\(\xi=c_b^{-1}\nabla b(U)\) is bounded, centered, has covariance
\(I\), and is invariant under signs and signed coordinate permutations.
It is affine on a fixed congruent triangulation of depth \(k_0\).
Take \(n=\lfloor m/(4k_0)\rfloor\) microsteps, evaluate the block
matrices \(V_j\) from completed driver sums, and consider the ideal
innovation martingale with increments \(V_j\xi_k/\sqrt n\).
Stop it after the first microstep whose state has norm at least
\(H+1\).  This is a predictable decision for subsequent increments;
the overshoot is \(O_V(n^{-1/2})\).

The invariance principle and continuity of the finitely many coefficient
maps give joint uniform-path convergence of the driver and the
unstopped martingale, block by block.  Their sphere-exit stopping
times and stopped paths converge too.  Indeed, within each fixed block
the limiting diffusion matrix is constant and uniformly invertible,
so its normal Brownian component crosses the sphere immediately upon
first hitting; hitting exactly at a deterministic block endpoint has
probability zero.  These facts give almost sure continuity of the
stopping map.  All stopped states lie in a fixed compact set, hence
their expected \(\Psi\) costs converge without a growth assumption.
The bounded reward sums converge to
\(\E\int_0^{1\wedge\tau_V}\sqrt{\det V_t}\,dt\).

\emph{The spatial branch law and its pathwise cost.}
At block \(j\), put \(a_j=\sqrt{\det V_j}\) and choose
\(A_j^{-T}=V_j/a_j\), so \(\det A_j=1\).  On an active triangle
\(Q=x_Q+d_QA_jO_QT\), add the signed bubble
\[
 b_Q(x)=\frac{d_Qa_j}{c_b\sqrt n}
 b\big(d_Q^{-1}O_Q^TA_j^{-1}(x-x_Q)\big),
 \qquad s_Q=\operatorname{sgn}W(b_Q).
 \tag{L8}\label{lower:eq:bubble}
\]
Its gradient increment is \(V_j\xi_Q/\sqrt n\), with
\(\xi_Q=s_Qc_b^{-1}O_Q\nabla b\). Before the first microstep, tile
the square in the initial chart \(A_1\) at mesh
\(2^{-\lfloor\sqrt m\rfloor}\), retaining complete triangles and
freezing the boundary remainder at gradient zero. Mark a parent
inactive after its state reaches the barrier; this requires no
retiling. At each later fixed control-block boundary retile the
surviving triangles in the new charts, at relative mesh
\(2^{-\lfloor\sqrt m\rfloor}\), retaining complete triangles and
discarding the boundary remainder. Compact
chart bounds give total discarded area \(o_V(1)\), and minimum
intrinsic leg \(d_Q\ge2^{-m/3}\), by the boundary-strip estimate for a triangle under two uniformly bounded charts.

Freeze \emph{any} assignment of all bubble signs and follow a uniform
spatial point.  Conditional on its parent and all previous child
choices, the point is uniform in that parent.  Congruent subdivision
therefore gives the same law of \(\xi_Q\) at every step: the fixed
sign and signed permutation do not change that law.  At a retiling,
conditional on the retained new triangle, uniformity holds again.
After geometric loss append independent artificial innovations.
After the state barrier do the same for the driver, leaving the
gradient fixed.  Induction proves that the ideal innovations are
independent with law \(\xi\), for every frozen sign assignment.

The actual spatial gradient is precisely this ideal state-stopped
martingale further stopped at geometric loss.  The filtration records
only parent and child identities and the predictable retiling decisions,
not the location inside a current triangle.  Thus conditional Jensen
at that stopping time gives
\[
 \int_D\Psi(\nabla u_m)
   \le\E_{\rm ideal}\Psi(\text{terminal stopped chain}).
 \tag{L9}\label{lower:eq:jensen}
\]
The right side is independent of all frozen Gaussian signs.  This is
a pathwise bound, not an estimate averaged over white noise.  It proves
\eqref{lower:eq:pathwise-cost}.  Discarded pieces freeze their old
gradient, so the asserted compact gradient bound holds everywhere.

\emph{Fresh Gaussian rewards.}
Reveal bubbles breadth first.  Every earlier bubble is disjoint from
\(Q\) or affine on \(Q\), so its pairing with \(b_Q\) vanishes by
the affine moments of \(b\).  Conditional on earlier Gaussian
observations, \(W(b_Q)/\norm{b_Q}_2\) is consequently a fresh standard
normal.  More explicitly, the directions take finitely many values at
each fixed depth; induction on this finite decision tree leaves an
isonormal Gaussian process on the orthogonal complement of the already
revealed directions.  For any past-measurable test \(v\), decomposition
into its \(b_Q\) component and the remaining Gaussian and past-measurable
components yields
\[
 \E[s_QW(v)\mid\text{past}]
 =\sqrt{2/\pi}\,\frac{\langle v,b_Q\rangle}{\norm{b_Q}_2}.
\]
In particular, with \(v=P_mb_Q\), the numerator is
\(\norm{P_mb_Q}_2^2\).  Poincar\'e on dyadic squares gives
\(\norm{P_mb_Q-b_Q}_2\le C_V2^{-m}d_Q^{-1}\norm{b_Q}_2=o(1)\norm{b_Q}_2\).
Determinant-one scaling in \eqref{lower:eq:bubble} gives
\(\norm{b_Q}_2=c_b' a_j|Q|/\sqrt n\), with universal \(c_b'>0\).
Thus each active parent has conditional expected reward at least
\(c a_j|Q|/\sqrt n\).  At each microstep their total coefficient is
the spatial expectation of \(a_j\) before the state barrier, up to
\(o_V(1)\) geometric loss.  Summing, using the stopped-chain convergence
and \(n\asymp m\), proves \eqref{lower:eq:spatial-reward}.
\end{proof}


% END bounded_bubble_realization.tex


% BEGIN general_pair_necessary.tex
\subsection{Necessary conditions for two arbitrary Young functions}

Let \(\theta_1,\theta_2\) be finite convex Young functions with
\(\theta_i(t)\ge t\), and put
\[
 r_i(S)=\frac{e^S}{\theta_i^{-1}(e^S)},\qquad S_j=2^j\quad(j\ge1).
 \tag{P1}\label{pair:eq:profiles}
\]
These functions are nondecreasing, since
\(\theta_i^{-1}(u)/u\) is nonincreasing.  The inverse-product inequality
proved above gives \(r_i(S)\asymp(\theta_i^*)^{-1}(e^S)\), with
absolute comparison constants.

\begin{lemma}[Finite weighted entropy]\label{pair:lem:finite-entropy}
Fix \(0<\gamma<e^{-1}\).  For finitely many positive weights \(v_j\), set
\[
 F(v)=\sup_{\substack{q_j\ge0\\\sum q_j\le\gamma}}
     \sum_jq_j\sqrt{\log(1/q_j)}\,v_j.
\]
Let \(T(v)>0\) be determined by
\(\sum_j\exp[-T(v)^2/v_j^2]=e^{-1}\).  Then
\[
 \gamma T(v)\le F(v)\le C_\gamma T(v),\qquad
 T(v)\ge\max_jv_j.
 \tag{P2}\label{pair:eq:entropy-comparison}
\]
Consequently, with \(v_j=\sqrt{S_j}/r_1(S_j)\), a uniform bound
\(F((v_j)_{j\le J})\le C r_2(S_J)\) is equivalent to the existence
of \(c>0\) such that, for every \(J\),
\[
 \sum_{j\le J}
 \exp\left[-c\frac{r_1(S_j)^2r_2(S_J)^2}{S_j}\right]\le e^{-1}.
 \tag{P3}\label{pair:eq:prefix12}
\]
\end{lemma}
\begin{proof}
The lower bound follows by taking
\(q_j=\gamma e\exp[-T^2/v_j^2]\): these probabilities sum to
\(\gamma\) and \(\log(1/q_j)\ge T^2/v_j^2\).  Also each term of
the defining sum for \(T\) is at most \(e^{-1}\), giving
\(T\ge\max v_j\).  For the upper bound put \(a_j=T/v_j\ge1\).
Split the normalized entropy sum according as
\(\sqrt{\log(1/q_j)}\le2a_j\) or not.  The first part is at most
\(2\gamma\), while the second is at most
\(C\sum_j e^{-2a_j^2}\le C e^{-1}\), exactly as in
Lemma~\ref{lower:lem:entropy}.  The last assertion follows from the
monotonicity in \(T\) of its defining exponential sum.  We will also use
\(T(v)\asymp\E\max_jv_j|G_j|\) for independent standard normals.
For the upper bound use the Gaussian union bound at levels \(sT\),
\(s\ge2\), and \(T^2/v_j^2\ge1\), then integrate the tail.
For the lower bound, the elementary bound
\(\PP(|G|\ge x)\ge c e^{-x^2}\), together with independence,
gives a fixed positive probability that the maximum exceeds \(T/2\).
\end{proof}

\begin{proposition}[Two necessary prefix conditions]
\label{pair:prop:necessary}
If the estimate in the question holds uniformly in the grid size and
sampling parameter, then \eqref{pair:eq:prefix12} holds for some
\(c>0\), and so does its axis-interchanged version
\[
 \sum_{j\le J}
 \exp\left[-c\frac{r_2(S_j)^2r_1(S_J)^2}{S_j}\right]\le e^{-1}
 \qquad(J\ge1).
 \tag{P4}\label{pair:eq:prefix21}
\]
The constant \(c\) may be enlarged to be the same in both conditions.
\end{proposition}
\begin{proof}
We prove the assertion by a bounded-state martingale construction.
The cost is now
\(\Psi(x,y)=\theta_1(|x|)+\theta_2(|y|)\).
The proof of Lemma~\ref{lower:lem:spatial} applies to this cost without
change: its cost argument uses only convexity and continuity on a
fixed compact set.  Thus arbitrarily large determinant reward at a
uniformly bounded expected terminal \(\Psi\) cost contradicts the
desired Gaussian, and hence empirical, estimate.

For clarity write
\[
 v_j=\frac{\sqrt{S_j}}{r_1(8S_j)},\qquad
 F_J=\sup_{\substack{q_j\ge0\\\sum_{j_0\le j\le J}q_j\le\gamma}}
       \sum_{j=j_0}^Jq_j\sqrt{\log(1/q_j)}\,v_j.
\]
We show that \(F_J/r_2(8S_J)\) must stay bounded.  Replacing
\(8S_j\) by \(S_j\) merely shifts indices by three and multiplies
the weights by \(\sqrt8\).  Omitting the fixed initial segment is
harmless: its entropy functional is finite, the functional is
subadditive over index sets, and \(r_2\ge1\).  Lemma
\ref{pair:lem:finite-entropy} then gives precisely
\eqref{pair:eq:prefix12}.

Choose probabilities within a factor two of the supremum defining
\(F_J\).  Their total mass can be taken to equal \(\gamma\), since
\(q\sqrt{\log(1/q)}\) is increasing for \(q\le\gamma\), with
\(\gamma\) chosen sufficiently small.  Put
\[
 y_0=\frac1{r_2(8S_J)},\quad H_j=e^{8S_j},\quad
 b_j=\frac{\theta_1^{-1}(H_j)}{y_0H_j}
     =\frac{v_j}{y_0\sqrt{S_j}},\quad
 \ell_j=q_jb_j,\quad K_j=\log(1+q_j^{-2}).
 \tag{P5}\label{pair:eq:pair-stage}
\]
Use the operational local-time stages of
Lemma~\ref{lower:lem:control}, now with
\(w_j(v)=S_j/[K_j(v+\ell_j^2)]\).  Their proofs were scale
invariant in \(b_j\), so they give
\[
 \PP(\text{failure in stage }j)\le Cq_j,\quad
 \E[R_j;\text{success}]\ge
       \frac{c q_j\sqrt{K_j}v_j}{y_0},\quad
 R_j\le\frac{2v_j}{y_0\sqrt{K_j}}.
 \tag{P6}\label{pair:eq:pair-resources}
\]
Condition on all the finitely many stages succeeding.  The resources
are independent and their sum has mean
\(\mu\ge cF_J/y_0\).  On the other hand
\[
 F_J\ge\gamma\sqrt{\log(1/\gamma)}\max_{j_0\le j\le J}v_j,
\]
by placing all the mass on a largest weight.  Hence the deterministic
maximum of the individual \(R_j\)'s is at most
\(C_\gamma F_J/y_0\le C_\gamma'\mu\).  Their conditional variance
is at most this maximum times \(\mu\).  Paley--Zygmund therefore
gives a constant \(p_\gamma>0\), independent of \(J\) and the two
profiles, such that with probability at least \(p_\gamma\) all
stages succeed and their total resource exceeds \(\mu/4\).
Here we used \(\PP(\text{all successes})\ge1-C\gamma\ge1/2\).

Set \(R=\mu/4\) and embed the stages in ordinary time by
\(dt=\sqrt{w_j(v)}\,dv/R\).  As before,
\(\sigma\beta=R\), \(dM=\sigma\,dB^1\), and the multiplicative
variance clock up to stage \(j\) is at most \(2S_j\).  This time start
the multiplicative martingale at \(y_0\):
\[
 Y_t=y_0\exp\left(\int_0^t\beta_r\,dB_r^2-s_t/2\right),
 \qquad X_t=Y_tM_t.
\]
Before the additional cap, \(\beta\) is still a function of \(B^1\)
alone.  Consequently the uncapped measure
\(dQ=(Y_1/y_0)d\PP\) preserves its entire \(B^1\) marginal.
Stop at the first crossing \(Y=y_0H_j\) during stage \(j\).
Under \(Q\), \(\log(Y/y_0)\) is Brownian motion plus half its
variance clock, conditionally on \(B^1\).  The cap probability is
therefore at most \(\sum_{j\ge j_0}2e^{-49S_j/4}\).
Choose \(j_0\) once so that this is less than \(p_\gamma/2\).
The stopped martingale \((X,Y-y_0)\) has triangular integrand
\(\left(\begin{smallmatrix}Y\sigma&MY\beta\\0&Y\beta\end{smallmatrix}\right)\)
before stopping.  Its expected determinant reward is at least
\[
 \frac{p_\gamma}{2}y_0\sqrt R
       \ge c_\gamma\sqrt{y_0F_J}.
 \tag{P7}\label{pair:eq:pair-reward}
\]

The expected terminal cost stays bounded, independently of \(J\).
Indeed, at a terminal state in stage \(j\), convexity and
\eqref{pair:eq:pair-stage} give
\[
 \theta_1(|X|)\le
 \theta_1\left(\frac{Y}{y_0H_j}\theta_1^{-1}(H_j)\right)
 \le Y/y_0.
\]
Also \(Y\le y_0H_J=\theta_2^{-1}(H_J)\), whence
\(\theta_2(Y)\le Y/y_0\).  Since the stopped \(Y\) is a bounded
martingale of mean \(y_0\le1\),
\[
 \E\{\theta_1(|X|)+\theta_2(|Y-y_0|)\}
 \le\E\theta_1(|X|)+\E\theta_2(Y)+\theta_2(y_0)
 \le2+\theta_2(1).
 \tag{P8}\label{pair:eq:pair-cost}
\]
The first inequality uses
\(|Y-y_0|\le\max(Y,y_0)\). Put \(C_0=3+\theta_2(1)\).
Apply the spatial lemma with a fixed error less than \(1/2\).
For all sufficiently fine meshes its pathwise energy is at most
\(C_0\), independently of \(J\). Dividing the potential by
\(C_0\) gives energy at most one by convexity, and divides its
reward by only this fixed constant. For every fixed finite list the state
and the integrand are bounded; in particular all compact-state
approximation and pathwise spatial cost arguments apply.  If
\(F_J/r_2(8S_J)=y_0F_J\) were unbounded,
\eqref{pair:eq:pair-reward}--\eqref{pair:eq:pair-cost} would contradict
the desired estimate.  This proves \eqref{pair:eq:prefix12}.
Interchanging the axes proves \eqref{pair:eq:prefix21}.
\end{proof}

\begin{remark}[A simpler pointwise necessary condition]
\label{pair:rem:pointwise}
The full uniform family of prefix sums implies
\[
 r_1(S)r_2(S)\ge c\sqrt S\qquad(S\text{ sufficiently large}).
 \tag{P9}\label{pair:eq:product}
\]
To see the needed uniform lower bound, put
$p_J=r_1(S_J)^2r_2(S_J)^2/S_J$. If $p_J$ is small, then for
$k\le J$ monotonicity gives
$r_1(S_k)^2r_2(S_J)^2/S_k\le2^{J-k}p_J$. Thus a block of
order $\log(1/p_J)$ preceding terms in the first prefix sum are bounded
below by $e^{-c}$, contradicting its uniform boundedness. There are
enough preceding terms because both inverse profiles have fixed positive
lower bounds. Dyadic values imply the full statement by monotonicity.
Equivalently,
\(\theta_1^{-1}(t)\theta_2^{-1}(t)
 \le C t^2/\sqrt{\log t}\).  The prefix conditions can be strictly
stronger: the examples below have both costs superlinear and satisfy
this pointwise condition while the desired estimate fails.
\end{remark}

\subsection{An exact subclass with comparable power profiles}

\begin{proposition}[Logarithmic power costs]\label{pair:prop:log-powers}
If $\theta_i(t)\asymp t(\log t)^{\beta_i}$ at infinity, with
$\beta_i\ge0$ and Young normalization near zero, the bound holds
exactly when
\[
 \beta_1+\beta_2>\tfrac12,
 \quad\text{or}\quad
 \beta_1+\beta_2=\tfrac12\ \text{and}\ \beta_1\beta_2>0.
\]
\end{proposition}
\begin{proof}
Here $r_i(S)\asymp S^{\beta_i}$. In the first orientation the
prefix exponents are comparable to
$2^{(2\beta_1-1)j+2\beta_2J}$.
If $\beta_1<1/2$ and the sum is at least $1/2$, they grow
geometrically as $J-j$ increases; at equality the rate is positive
exactly when $\beta_2>0$. If $\beta_1=1/2$ and $\beta_2>0$,
the sum is at most $J\exp(-c2^{2\beta_2J})$.
If $\beta_1>1/2$, the sum is bounded by the convergent infinite
series $\sum_j\exp(-c2^{(2\beta_1-1)j})$.
Apply the same argument after interchanging coordinates.
If the sum of exponents is smaller than $1/2$, the pointwise product
condition fails. At $(\beta_1,\beta_2)=(1/2,0)$ the first entropy
weights are constant, and their maximum grows as $\sqrt{\log J}$,
so the corresponding prefix condition fails. The other endpoint is
identical with the coordinates interchanged.
\end{proof}

\begin{proposition}\label{pair:prop:power-comparable}
Suppose that for some fixed \(\kappa>0\),
\(r_1(S)\asymp r_2(S)^\kappa\) for all sufficiently large \(S\).
Then the desired estimate holds if and only if
\eqref{pair:eq:product} holds.  In particular, if
\(\theta_1=\theta_2=\theta\), its exact condition is
\[
 \theta(t)\ge c\,t(\log t)^{1/4}
 \qquad(t\text{ sufficiently large}).
 \tag{P10}\label{pair:eq:symmetric}
\]
\end{proposition}
\begin{proof}
Necessity follows from Proposition~\ref{pair:prop:necessary}.
Conversely, power comparability and \eqref{pair:eq:product} imply
\[
 r_1(S)\ge cS^{\beta_1},\quad
 r_2(S)\ge cS^{\beta_2},\qquad
 \beta_1=\frac{\kappa}{2(\kappa+1)},\quad
 \beta_2=\frac1{2(\kappa+1)}.
\]
Both exponents are positive and sum to \(1/2\).
These lower bounds imply the two prefix conditions directly: the
first normalized inverse product is at least
$c2^{\beta_2(J-j)}$ and its axis-swapped counterpart is at least
$c2^{\beta_1(J-j)}$. Both exponential sums are bounded geometric-tail
sums. The general prefix criterion therefore proves sufficiency.
Equivalently the gradient costs dominate
\(t(\log t)^{\beta_i}\).

When the costs are equal, \(\kappa=1\) and the inverse condition is
\(r(S)\ge cS^{1/4}\).  For any fixed \(\beta>0\),
\(r(S)\ge cS^\beta\) is equivalent to
\(\theta(t)\ge c't(\log t)^\beta\) at infinity.  For one direction
put \(S=\log\theta(t)\ge\log t\).  For the other put
\(t=\theta^{-1}(e^S)\): if \(t\le e^{S/2}\), then
\(e^S/t\ge e^{S/2}\); otherwise \(\log t\ge S/2\), and the stated
cost lower bound gives \(e^S/t\ge c'(S/2)^\beta\).
This proves \eqref{pair:eq:symmetric}.
\end{proof}

\subsection{An iterated-logarithm family near the hard endpoint}

The following calculation separates the pointwise product obstruction
from the extra entropy obstruction.  All tail comparisons are for
sufficiently large arguments; functions may be patched near zero to be
Young functions without changing the conclusions.

\begin{proposition}\label{pair:prop:iterated}
Suppose
\begin{align*}
 \theta_1(t)&\asymp
 t(\log t)^{1/2}(\log\log t)^{\gamma_1}
                    (\log\log\log t)^{\delta_1},\\
 \theta_2(t)&\asymp
 t(\log\log t)^{\gamma_2}(\log\log\log t)^{\delta_2},
\end{align*}
where either \(\gamma_2>0\), or \(\gamma_2=0\) and
\(\delta_2\ge0\), as required for the second cost to dominate
\(t\).  Then the desired estimate holds exactly when
\[
 \gamma_1+\gamma_2>0,
 \quad\hbox{or}\quad
 \gamma_1+\gamma_2=0\ \hbox{ and }\ \delta_1+\delta_2\ge\tfrac12.
 \tag{P11}\label{pair:eq:iterated-condition}
\]
The two prefix conditions, and hence the general theorem, give this
necessary and sufficient condition.
In particular, for every \(a>0\), the two superlinear costs
\[
 \theta_1(t)\asymp\frac{t\sqrt{\log t}}{(\log\log t)^a},
 \qquad \theta_2(t)\asymp t(\log\log t)^a
\]
fail the desired estimate, although their product has the pointwise
critical order \(t^2\sqrt{\log t}\).
\end{proposition}
\begin{proof}
Inversion of these iterated-logarithm functions gives
\[
 r_1(S_j)\asymp2^{j/2}j^{\gamma_1}(\log j)^{\delta_1},\qquad
 r_2(S_j)\asymp j^{\gamma_2}(\log j)^{\delta_2}.
\]
Thus the first entropy weights are
\(v_j\asymp j^{-\gamma_1}(\log j)^{-\delta_1}\).
For independent standard normals \(G_j\), their entropy functional
is comparable to \(\E\max_{j\le J}v_j|G_j|\), or directly to
the exponential-sum threshold in
Lemma~\ref{pair:lem:finite-entropy}.  The following elementary
estimates therefore suffice:
\[
 F_J\asymp
 \begin{cases}
 J^{-\gamma_1}(\log J)^{1/2-\delta_1},&\gamma_1<0,\\
 1+(\log J)^{\max(1/2-\delta_1,0)},&\gamma_1=0,\\
 1,&\gamma_1>0.
 \end{cases}
 \tag{P12}\label{pair:eq:iterated-entropy}
\]
For increasing polynomial weights the last \(J/2\) indices give the
lower bound, and the usual Gaussian maximum bound gives the upper
bound.  Polynomially decaying weights have a summable inverse-square
exponential series.  For pure logarithmic weights with
\(\delta_1<1/2\), put
\(T=C(\log J)^{1/2-\delta_1}\); for \(j\le J\),
\(T^2(\log j)^{2\delta_1}\ge C^2\log j\) when
\(0\le\delta_1<1/2\), so the exponential-sum upper bound follows.
Negative \(\delta_1\) is covered by the increasing-weight argument.
If \(\delta_1\ge1/2\), the infinite exponential series converges
for a sufficiently large constant.  These observations prove
\eqref{pair:eq:iterated-entropy}.

The swapped weights are
\(\sqrt{S_j}/r_2(S_j)\asymp
2^{j/2}j^{-\gamma_2}(\log j)^{-\delta_2}\).
Since they grow geometrically, their entropy functional is comparable
to the final weight.  Consequently the two necessary prefix bounds
are equivalent, in this family, to \eqref{pair:eq:iterated-condition}.
Proposition~\ref{pair:prop:necessary} proves failure in every excluded
case.

Conversely the same entropy calculation verifies both prefix conditions
in every included case. The general two-function criterion proves the
matching upper bound, including equality in
\eqref{pair:eq:iterated-condition}.
\end{proof}

\begin{remark}
The examples above are consequences of the general two-function
criterion; they impose no additional regularity beyond the finite Young
function hypotheses. The explicit iterated-logarithm comparisons only
make that criterion easier to read in this subclass.
\end{remark}

% END general_pair_necessary.tex


\section{Completion of the Young-function theorem}
The prefix conditions give the allocated phase and the exact Gaussian
current of Proposition~\ref{pair:prop:flow}. For every potential $f$
of continuum modular energy at most one, weak integration by parts
and the complementary Young inequality yield
\[
 |W_h^0(f)|\le
 2\sum_{i=1}^2\|J_i\|_{L^{\psi_i}}\|\partial_i f\|_{L^{\theta_i}}
 \le4\max_i\|J_i\|_{L^{\psi_i}}.
\]
The grid--continuum comparison therefore proves the Gaussian upper
bound. Conversely, if either prefix condition fails, the bounded-cost
local-time controls and bounded-state bubble realization give
deterministic-energy potentials with arbitrarily large multiples of
$\sqrt p$ as expected pairing. Cell averaging transfers these to the
lattice without additional energy. This proves the converse Gaussian
implication. The upper empirical comparison and the uniform-law
Skellam converse prove the equivalence with (17.18), including
$m_0\asymp p$. The finitely many coarse meshes omitted during the
constructions are absorbed by enlarging the fixed constant. This proves
Theorem~\ref{pair:thm:criterion}.\qed


% BEGIN finite_Borel_relaxation.tex
\section{Finite Borel costs and quantitative convexification}
\label{sec:finite-borel-relaxation}

For $j=1,2$, let $\theta_j:[0,\infty)\to[0,\infty)$ be finite Borel
functions satisfying $\theta_j(t)\ge t$. Let $\phi_j$ be the even
convex envelope of $s\mapsto\theta_j(|s|)$, and put
\[
 a_j=\inf\theta_j=\phi_j(0),\quad a=a_1+a_2,\qquad
 \vartheta_j(t)=\phi_j(t)-a_j+t\quad(t\ge0).
\]
The equality at zero follows by mixing $s$ and $-s$. Since
$\phi_j(s)\ge|s|$, the functions $\vartheta_j$ are finite convex
Young functions, vanish at zero, and dominate $t$.

\begin{theorem}\label{thm:finite-borel-pair-relaxation}
If $a<1$, the original pair satisfies the discrepancy bound if and
only if the normalized Young pair $(\vartheta_1,\vartheta_2)$
satisfies the criterion proved above. If $a\ge1$, the original
pair always satisfies the discrepancy bound.
\end{theorem}

We first give the elementary finite local lemma needed to handle
costs that are unbounded on every neighborhood.

\begin{lemma}[Finite local lemma]\label{lem:finite-local-lemma}
Suppose $E_1,\ldots,E_N$ are events determined by independent
variables, and each event shares variables with at most $d$ others.
If, for some $0<x<1$,
$\mathbb P(E_i)\le x(1-x)^d$ for every $i$, then
\[
 \mathbb P\Big(\bigcap_i E_i^c\Big)\ge(1-x)^N>0.
\]
\end{lemma}

\begin{proof}
Inductively prove that, for every $S$ not containing $i$,
$\mathbb P(E_i\mid\bigcap_{j\in S}E_j^c)\le x$.
Partition $S$ into the neighbors $A$ of $i$ and the remaining
indices $B$. The event $E_i$ is independent of the entire family
indexed by $B$. Its conditional probability after also excluding
the events in $A$ is therefore at most
\[
 \frac{\mathbb P(E_i)}
 {\mathbb P(\bigcap_{j\in A}E_j^c\mid\bigcap_{j\in B}E_j^c)}
 \le \frac{x(1-x)^d}{(1-x)^{|A|}}\le x.
\]
The middle bound follows by the induction hypothesis and the
chain rule; its conditioning sets have smaller cardinality than
$S$. The same chain rule now gives the asserted lower bound.
\end{proof}

\begin{proof}[Proof of Theorem~\ref{thm:finite-borel-pair-relaxation}]
For sufficiency when $a<1$, note that
$\vartheta_j(t)\le\phi_j(t)+t\le2\theta_j(t)$.
Dividing an original admissible function by two puts it in the
normalized Young ball. We prove the converse by transferring the
bounded-gradient Gaussian lower tests from the Young-pair theorem.

Fix arbitrarily large $D$. That lower construction, followed by a
fixed rescaling, provides $H<\infty$ and, for arbitrarily large $m$,
random functions $u_m$ on the unit square with
\[
 \int u_m=0,\quad |\partial_j u_m|\le H,\quad
 \int\sum_j\vartheta_j(|\partial_j u_m|)\le\eta,
 \qquad
 \mathbb E\langle P_{2^{-m}}W^0,u_m\rangle\ge D\sqrt m.
\]
Here $\eta>0$ is fixed and sufficiently small relative to $1-a$.
Their convex-envelope energy is at most $a+\eta$.
Conditional expectation given the fixed coarse histogram preserves
these pathwise convex energy and gradient bounds and preserves the
coarse reward. Indeed, take conditional expectations on a countable
dense set of spatial points; the uniform Lipschitz bound gives a
continuous version. Fubini in the weak derivative identity identifies
its derivatives with the conditional expectations of the original
derivatives, and Jensen gives the pathwise energy bound. Hence assume
$u_m$ is measurable with respect to
$P_{2^{-m}}W^0$.

\paragraph{Finite base gradients and finite well slopes.}
Extend $u_m$ evenly by reflection to the period-two torus and
convolve with a coordinatewise-even nonnegative probability kernel at
scale $\rho=2^{-2m}$.
The result $u_m^\rho$ has coordinate gradients bounded by $H$,
Hessian bounded by $CH/\rho$, and uniform error at most $CH\rho$.
Jensen and reflection symmetry show that its convex-envelope
energy on the square has not increased.

Triangulate the square by a grid of side $\delta=2^{-3m}$ and a
fixed right-triangle diagonal. Round the values of $u_m^\rho$ at
vertices to the nearest multiple of $q\delta$, and interpolate
affinely to obtain $\widehat u_m$. Here $q>0$ is fixed and small.
The interpolation estimates are
\[
 \|\widehat u_m-u_m^\rho\|_\infty\le C(H+q)\delta,
 \qquad
 |\nabla\widehat u_m-\nabla u_m^\rho|
                  \le CH\delta/\rho+Cq.
\]
Its gradients therefore lie in a fixed finite subset of
$q\mathbb Z^2$, independent of $m$. Continuity of the convex
$\phi_j$ on compact intervals permits choosing $q$, and then $m$,
so that its convex-envelope energy is at most $a+2\eta$.
To justify the finite-well assertion, the infimum of the cost of all
finite convex combinations with prescribed mean is a finite convex
minorant, hence equals $\phi_j$. In any one approximating finite
combination, minimizing the linear cost under its mass and mean
constraints has an extreme solution supported at at most two points.
For each coordinate of each base gradient, choose such a two-point
original-well mixture with that mean and cost at most the
corresponding $\phi_j$ value plus $\eta$. Only finitely many well
slopes are needed, and their original costs are finite.

\paragraph{Two clipped waves.}
On a base triangle with gradient $(b_1,b_2)$, add a nonnegative
periodic triangular wave in $x$, of amplitude at most
$\epsilon_1=2^{-4m}$, whose derivative changes $b_1$ to its two
chosen first-coordinate wells in the prescribed proportions.
Clip this wave by the distance to the triangle boundary. The
correction vanishes at the boundary, is continuous and piecewise
affine, and has bounded gradients. Outside a boundary strip of
width $C\epsilon_1$, its first derivative is an original well and
its second derivative remains $b_2$.

On every retained affine polygon, add a corresponding triangular
wave in $y$, of amplitude at most $\epsilon_2=2^{-5m}$, clipped by
distance to that polygon's boundary. Keep the first boundary strips
unchanged. Denote the resulting continuous piecewise affine function
by $v_m$. It agrees with $\widehat u_m$ on all base triangle boundaries,
has gradient bounded by a fixed $K$, and
\[
 \|v_m-\widehat u_m\|_\infty\le\epsilon_1+\epsilon_2.
\]
Call the interiors where both derivatives are original well slopes
the good regions. The remaining region $B_m$ satisfies
\[
 |B_m|\le C(\epsilon_1/\delta+\epsilon_2/\epsilon_1).
\]
Indeed the total perimeter of the original triangles is
$O(\delta^{-1})$ and that of the first retained polygons is
$O(\epsilon_1^{-1})$. A polygon's inner boundary strip has area
at most its perimeter times its width. The total interface length
of the final affine partition is at most $C/\epsilon_2$, and its
number of vertices is at most $C/(\epsilon_1\epsilon_2)$.
These bounds follow by slicing the original triangles into their
first-wave stripes and each resulting polygon into second-wave
stripes; clipping adds a bounded number of segments per stripe.
More explicitly, there are $O((\delta\epsilon_1)^{-1})$ first
polygons. Each has height at most $C\delta$; the second slicing
therefore gives $O((\epsilon_1\epsilon_2)^{-1})$ pieces in total.
All constants are uniform because the base gradients and well slopes
come from fixed finite libraries.

Only well costs are needed on the good regions. Their summed
integral is at most
\[
 a+4\eta+C_\mathrm{well}
       (\epsilon_1/\delta+\epsilon_2/\epsilon_1).
\]
For clarity, the well proportions on a polygon differ from their
prescribed means by at most the wave period times its perimeter:
integrate by parts a bounded periodic primitive of the zero-mean
stripe indicator. Summing the first and second wave errors gives
the displayed bound. Removing the boundary strips cannot increase
these nonnegative well costs. We make no assertion about the
original cost of the continuum gradients in $B_m$.

\paragraph{A bounded-cost assignment on affected lattice edges.}
Take $h=2^{-8m}$ and initially sample $v_m$ at the centers of the
$h$-squares. Declare a grid edge defective if its joining segment
is not contained in a single good affine region. Mark both endpoints
of every defective edge, and call all edges incident to a marked
vertex affected. The normalized number of affected edges is at most
\[
 C\bigl(|B_m|+h/\epsilon_2+h^2/(\epsilon_1\epsilon_2)\bigr)=o(1).
\]
This is the area count for $B_m$, plus the standard grid-edge
count along a polygonal interface; marking endpoints multiplies
the count by only a numerical constant.
To verify that count directly, cover each interface segment by a
rectangle of width $Ch$ and cover each endpoint by a disk of radius
$Ch$. Their total area is at most $Ch$ times the interface length
plus $Ch^2$ times its number of vertices. Every defective grid edge
lies in $B_m$ or in this cover, and disjoint grid squares convert
this area bound to the displayed normalized count.

Fix $\tau=1/10$ and perturb each marked vertex value independently
by $h\xi_v$, with $\xi_v$ uniform on $[-\tau,\tau]$; keep other
vertices fixed. Every affected edge has at least one random endpoint.
Its derivative equals a fixed number in $[-K,K]$ plus the difference
of at most two such uniforms, so it has density at most $1/(2\tau)$
on $[-K-2\tau,K+2\tau]$. Since the original costs are finite Borel
functions, their superlevel sets on this compact interval have
Lebesgue measures tending to zero. Choose one fixed $M$ large enough
that, for each direction, the probability its affected-edge cost
exceeds $M$ is at most $1/100$.

Each bad-cost event depends on at most the two endpoint jitters and
shares variables with at most six other edge events. Since
$1/100<(1/8)(7/8)^6$, Lemma~\ref{lem:finite-local-lemma} supplies
an assignment for which every affected-edge cost is at most $M$.
No assignment can change any value by more than $\tau h$.
For fixed $m$ there are only finitely many rounded base vertex arrays,
because $|u_m|\le2H$. Choose a successful assignment once for each
such array. The resulting assignment rule is measurable and depends
only on the coarse histogram. This avoids any measurable-selection
or algorithmic assumption.

Let $z_m$ be the cellwise constant function given by these perturbed
samples. Unaffected edges still see exact original-well derivatives.
Their area-counting discrepancy is controlled by the same interface
count above, and affected edges cost at most $M$. Consequently its
normalized original lattice energy is at most
\[
 a+4\eta+o(1).
\]
Choose $\eta$ with $a+5\eta<1$. For all sufficiently large $m$,
$z_m$ is original-cost admissible. Neither integrability nor local
boundedness of the original costs was used.

The same proof also permits either original cost to be locally bounded
without being measurable. On the fixed compact interval of affected
slopes choose $M$ above that cost's supremum. Its bad-edge events are
then empty, so no measurability is needed for that coordinate. The other
coordinate may independently be measurable or locally bounded. In
particular this proves the stated classification for every locally
bounded pair, with no regularity assumption beyond local boundedness.

\paragraph{The Gaussian reward and depth.}
The function $z_m$ is coarse-histogram measurable. Added fine
Gaussian bridges therefore have conditional mean zero, giving
\[
 \mathbb E\langle P_hW^0,z_m\rangle
       =\mathbb E\langle P_{2^{-m}}W^0,z_m\rangle.
\]
Furthermore
\[
 \|z_m-u_m\|_\infty
 \le C(H+q)\rho+\epsilon_1+\epsilon_2+(K+\tau)h,
 \qquad \mathbb E\|P_{2^{-m}}W^0\|_1\le C2^m.
\]
The expected reward is thus $D\sqrt m-o(1)$. Dividing the fine cell
values by $h$ gives a lattice function with budget $h^{-2}$ and
standard lattice Gaussian reward $h^{-2}(D\sqrt m-o(1))$.
The new depth is $p=8m$, a fixed factor. Since $D$ is arbitrary,
the Gaussian estimate fails for the original pair, and the
Gaussian--empirical lower comparison transfers this failure to the
sampling regime of the book.

\paragraph{Baseline degeneracies.}
If $a>1$, the baseline $a n(n-1)$ exceeds $n^2$ for all sufficiently
large $n$; only finitely many sizes remain. Their bounds follow by
finite dimensionality and the empirical--Gaussian upper comparison.
If $a=1$, admissibility and $\theta_j(t)\ge\max(a_j,t)$ imply
\[
 \sum_{j,e}(|D_jh|-a_j)_+\le n.
\]
The linear dyadic Gaussian divergence inverse $J=TG^0$ satisfies
$\mathbb E\sum_e|J_e|\le Cn^2\sqrt p$ and
$\mathbb E\|J\|_\infty\le Cp$. Summation by parts gives
\[
 |\langle G^0,h\rangle|
 \le\sum_j a_j\sum_{e\parallel j}|J_e|+n\|J\|_\infty.
\]
Its expectation is at most $Cn^2\sqrt p+Cnp\le Cn^2\sqrt p$.
The empirical upper comparison proves the asserted estimate.
\end{proof}

% END finite_Borel_relaxation.tex


\section{Why completely nonmeasurable costs are different}

% BEGIN pathological_pair_counterexample.tex
\subsection{A genuine obstruction to unrestricted two-cost convexification}
\label{sec:hamel-pair-pathology}

For two completely arbitrary finite functions, replacing the costs by
their even convex envelopes is not a valid general principle. The
following example uses functions that are not measurable; the literal
formulation $\theta_j(t)\ge t$ does not exclude them.

\begin{proposition}\label{prop:hamel-pair-pathology}
There are finite functions $\theta_1,\theta_2:[0,\infty)\to[0,\infty)$
with $\theta_j(0)=0$ and $\theta_j(t)\ge t$ such that:
\begin{enumerate}
\item both even convex envelopes are exactly $t\mapsto|t|$;
\item the desired $O(n^2\sqrt{pm_0})$ empirical discrepancy estimate
holds for the pair $(\theta_1,\theta_2)$.
\end{enumerate}
Thus the unrestricted pair problem cannot be classified solely by
these pointwise convex envelopes, since the linear--linear pair has
already been shown to fail the estimate.
\end{proposition}

\begin{proof}
Choose a Hamel basis of $\mathbb R$ over $\mathbb Q$ containing
$1$ and $\sqrt2$. Let $P:\mathbb R\to\mathbb Q$ extract the
coefficient of $1$. It is additive, $P^2=P$, and both its image
and kernel are dense: they contain $\mathbb Q$ and
$\sqrt2\mathbb Q$, respectively. Set, for $t\ge0$,
\[
 \theta_1(t)=t+|P(t)|^2,\qquad
 \theta_2(t)=t+|t-P(t)|^2.
\]
These functions are finite and dominate $t$. The first equals $t$
on a dense subset, as does the second. Every finite convex minorant
of either even cost is continuous, hence is at most $|t|$ by density;
$|t|$ itself is a minorant. This proves the first assertion.

Let $h$ satisfy the original budget $n^2$, and write
$a(v)=P(h(v))$, $b(v)=h(v)-a(v)$. Additivity of $P$, and its oddness,
give the exact identity
\[
 \operatorname{TV}(h)
   +\sum|D_xa|^2+\sum|D_yb|^2\le n^2.
\]
There are fewer than $n^2$ edges in each direction. Cauchy--Schwarz
therefore gives
\[
 \sum|D_xa|\le n^2,\qquad \sum|D_yb|\le n^2,
\]
and the identity $h=a+b$ implies
\[
 \sum|D_ya|\le\sum|D_yh|+\sum|D_yb|\le2n^2,
 \qquad
 \sum|D_xb|\le2n^2.
\]
Consequently $a$ lies in four times the anisotropic Young ball with
costs $(t+t^2,t)$, while $b$ lies in four times the ball with costs
$(t,t+t^2)$. Both balls satisfy the desired estimate by
Theorem~\ref{pair:thm:criterion}: the linear complementary width is
constant, whereas the quadratic-cost complementary width grows
exponentially in the log-level. The resulting prefix sums are bounded.

The centered empirical functional is linear in the test function,
so its value at $h=a+b$ is bounded by the sum of its absolute values
at $a$ and $b$. Taking the two suprema and expectations proves the
second assertion, with a universal constant. The same argument
works directly for Gaussian width.
\end{proof}

Finite Borel costs, including those unbounded on compact sets, are
covered by the preceding quantitative relaxation theorem. This example
lies outside that measurable setting and shows that the qualification
cannot simply be deleted from a convex-envelope classification.

% END pathological_pair_counterexample.tex


\paragraph{Provenance of the methods.}
The discrepancy problem and its normalization are from Talagrand's
second edition, Chapter 17. The local Gaussian current and adaptive
bubble constructions were developed in the Ultimatematching project,
in particular the Goldman--Grotto--Trevisan source files
\texttt{umt.tex} and \texttt{endpoints.tex}. All ingredients used here
are proved in this note. The two prefix conditions, width allocation,
rounded phase, two-cost repair estimate, and measurable-cost lattice
relaxation are the extensions developed in this review.

\begin{thebibliography}{9}
\bibitem{talagrand}
M. Talagrand, \emph{Upper and Lower Bounds for Stochastic Processes:
Decomposition Theorems}, second edition, Springer, 2021,
Research Problem 17.2.2.
\bibitem{umt}
M. Goldman, F. Grotto, and D. Trevisan,
\emph{The Ultimate Matching Theorem}, preprint, 2026.
\url{https://hexagonmath.org/2611.00294v1}.
\bibitem{repository}
Ultimatematching project source files,
\url{https://github.com/grottofrancesco/UltimateMatchingTheorem}.
\end{thebibliography}
\end{document}
